%%%BLOCK id=010-03 ch=04 sec=竖直面圆周·绳球 kind=method alt=06 cont=none y=0.25-0.43
\textbf{绳球}\quad $T\perp V$，$W=0$

上到最高点\ $V\ge\sqrt{gr}$

%FIG p010_a 0.50 0.282 0.585 0.342
\notefig[0.25]{figures/p010_a.png}

上不去\ $\left\{\begin{array}{ll}\text{上半圆} & \text{斜上抛}\\ \text{下半圆} & \text{单摆}\end{array}\right.$

类绳模型$\to$弹力只能向内

最高点最低点拉力差\ $6mg$
%%%END
%%%BLOCK id=010-04 ch=04 sec=竖直面圆周·杆球 kind=method alt=06 cont=none y=0.44-0.68
\textbf{杆球}\quad 杆力\ $W_{\text{总}}=0$

上到最高点\ $V_{\text{高}}\ge0$

上不去\quad 原路反回 % “反回”疑为“返回”

类杆模型\quad 小球受弹力可内可外

%FIG p010_b 0.525 0.443 0.975 0.602
\notefig[0.6]{figures/p010_b.png}
% 图 p010_b：一条大波浪形轨道，沿途分段标注“杆”“绳”“杆”

%FIG p010_c 0.412 0.576 0.66 0.672
\notefig[0.35]{figures/p010_c.png}
% 图 p010_c：四个小示意图（圆环形“杆”、拱形“杆”、碗形“绳”、圆形“杆”）
%%%END
%%%BLOCK id=041-01 ch=04 sec=关联速度 kind=method alt=02 cont=none y=0.03-0.25
\textbf{运动关联}% 页面左上角标题(带下划线框线)；页左侧另有一条红色大括号，自本块“人 V”图起括至“同轴同ω”图止

\noindent $V_{\text{船}}\cos\alpha=V_{\text{人}}$\quad 船\unclear{比}人快% 页面右上角两行(第二行中间一字辨认不清)

\begin{itemize}
\item 人匀速船加速
\item 人：减速船才有可能匀速
\end{itemize}

%FIG p041_a 0.07 0.125 0.30 0.208
\notefig[0.28]{figures/p041_a.png}
% 图内标注：人 $V$(岸上，水平绳段)；绳方向分速度 $V_{\parallel}$ 改变绳长；船速 $V_1$、夹角 $\alpha$；垂直分速度 $V_{\perp}$(斜向下箭头)
\noindent $V_{\perp}$ 绳子摆动(圆周运动)% 图正下方的标注，因紧邻右侧文字未含在图内

\noindent $V=V_1\cos\alpha$\qquad 只能分解实际速度

%FIG p041_b 0.63 0.15 0.705 0.245
\notefig[0.25]{figures/p041_b.png}
% 图内标注：船受力图——$F$(沿绳斜向上，带虚线分量)、$mg$ 向下
\noindent $F\sin\alpha=mg$
%%%END
%%%BLOCK id=041-02 ch=04 sec=关联速度 kind=method alt=none cont=none y=0.25-0.34
%FIG p041_c 0.115 0.25 0.36 0.315
\notefig[0.28]{figures/p041_c.png}
% 图内标注：绳绕过上方定滑轮连接左右两物体；左物体速度 $V_1$(水平向右)，绳与水平夹角 $\alpha$；右物体速度 $V_2$(水平向右)，绳与水平夹角 $\beta$；虚线箭头、沿绳/垂直绳的分量箭头
\[
V_1\cos\alpha=V_2\cos\beta
\]
%%%END
%%%BLOCK id=041-03 ch=04 sec=关联速度 kind=method alt=03 cont=none y=0.24-0.35
%FIG p041_d 0.42 0.235 0.605 0.335
\notefig[0.25]{figures/p041_d.png}
% 图内标注：水平面上物体 $B$ 向左运动(标 $V\cos\alpha$ 的分量)，绳绕过上方定滑轮另一端竖直挂物体 $A$(向上箭头)，滑轮处标 $V_A$，夹角 $\alpha$
\noindent $V_A=V\cos\alpha$\quad $\alpha\downarrow$\quad $\cos\alpha\uparrow$

\noindent $V_A\uparrow$\quad $\therefore T>mg$
%%%END
%%%BLOCK id=041-04 ch=04 sec=关联速度 kind=method alt=none cont=none y=0.35-0.45
%FIG p041_e 0.12 0.355 0.285 0.445
\notefig[0.25]{figures/p041_e.png}
% 图内标注：直角三角形(斜边为杆/绳)，上端点 1 处标 $V_1$ 竖直向下及夹角 $\theta$，下端点 2 处标 $V_2$ 水平向右、夹角 $\theta$、虚线箭头
\[
V_1\cos\theta=V_2\sin\theta
\]
%%%END
%%%BLOCK id=041-05 ch=04 sec=关联速度 kind=note alt=none cont=none y=0.45-0.49
\red{\struck{\unclear{难}}\ 沿绳方向速度相等}% 红笔；最左一字被红笔大叉划去(像“难”)；此红字是对上方各例的总结
%%%END
%%%BLOCK id=041-06 ch=04 sec=关联速度 kind=method alt=03 cont=none y=0.46-0.61
%FIG p041_f 0.11 0.46 0.30 0.61
\notefig[0.25]{figures/p041_f.png}
% 图内标注：定滑轮(顶部横梁)，绳一端竖直挂物体 2(向上箭头)，另一端斜拉竖直轨道内的物体 1；物体 1 处标夹角 $\theta$，粗箭头向下标 $V_1$
\noindent $V_1\cos\theta=V_2$

\noindent $\theta\downarrow$\quad $V_1\cos\theta\uparrow$\quad $\therefore V_2\uparrow$\quad 2向上加速

\noindent $T>m_2g$
%%%END
%%%BLOCK id=041-07 ch=04 sec=关联速度 kind=method alt=none cont=none y=0.60-0.70
%FIG p041_g 0.07 0.605 0.268 0.685
\notefig[0.27]{figures/p041_g.png}
% 图内标注：斜杆(长 $L$，下端 $B$ 铰在斜面/地面左下，与水平夹角 $\theta$)，杆上端点 $A$；右下方块(高 $h$)以速度 $V$ 向右，其上角与杆接触，红箭头为垂直杆的分速度；红笔沿杆描出
\noindent \fcolorbox{red!75!black}{white}{\red{同轴同$\omega$}}% 红笔圈出，由 $A$ 处括号引出

\[
\left\{\begin{aligned}
&V\sin\theta=\omega\,\dfrac{h}{\text{\struck{\unclear{$\tan\theta$}}}\,\sin\theta}\\
&V_A=\omega L
\end{aligned}\right.
\]% 分母中最左一项被涂划(像 tanθ)，其后为 sinθ；原稿此处 ω 写得像 W

\noindent \red{同一转轴 $\omega$ 相同}
%%%END
%%%BLOCK id=041-08 ch=04 sec=关联速度 kind=method alt=none cont=none y=0.71-0.82
%FIG p041_h 0.07 0.715 0.245 0.81
\notefig[0.25]{figures/p041_h.png}
% 图内标注：一个五边形物体速度 $V_2$ 由绳(与水平夹角 $\alpha$)拉着；绳另一端为人，人速度 $V_1$ 水平向右(红箭头为沿绳/垂直绳分量)，虚线沿绳延长
\[
V_2=V_1\cos\alpha
\]
%%%END
%%%BLOCK id=042-01 ch=04 sec=关联速度·交点速度 kind=derivation alt=none cont=none y=0.02-0.12
\textbf{交点速度}% 页面左上角标题(后接虚线)

%FIG p042_a 0.12 0.027 0.34 0.112
\notefig[0.3]{figures/p042_a.png}
% 图内标注：两条直线(其中一条水平线上标 $V_2\Delta t/\sin\alpha$)交于 $M$，$V_2$(小字)，$V_1\Delta t/\sin\alpha$、$V_1\Delta t$、向下箭头 $V_1$；$M$ 处向下箭头 $V_M$；虚线斜线末端箭头 $V_2$
%FIG p042_b 0.40 0.032 0.545 0.095
\notefig[0.25]{figures/p042_b.png}
% 图内标注：三个矢量——向左箭头标 $V_2\Delta t/\sin\alpha$，向左下箭头标 $V_M$，向右下箭头标 $V_1\Delta t/\sin\alpha$
\[
V_M=\frac{\sqrt{\left(\dfrac{V_2\Delta t}{\sin\alpha}\right)^{2}+\left(\dfrac{V_1\Delta t}{\sin\alpha}\right)^{2}-2\left(\dfrac{V_1\Delta t}{\sin\alpha}\,\dfrac{V_2\Delta t}{\sin\alpha}\right)}}{\Delta t}
\]% CHECK: 余弦定理的交叉项按理应含 cos(夹角)，原稿括号内未见 cos 因子，照原样转录
%%%END
%%%BLOCK id=042-02 ch=04 sec=关联速度·交点速度 kind=derivation alt=none cont=none y=0.12-0.25
%FIG p042_c 0.05 0.12 0.28 0.21
\notefig[0.25]{figures/p042_c.png}
% 图内标注：两个相互交叠的菱形(旋转45°的方块)，左边菱形上标 $V$(向左箭头)，右边菱形上标 $2V$(向右箭头)，交点标 $M$
小量分析，延长\ 找角。

%FIG p042_d 0.33 0.152 0.568 0.256
\notefig[0.28]{figures/p042_d.png}
% 图内标注：两条45°交叉直线交于 $M$(沿一线标 $\sqrt2V\Delta t$)，左下角 $45^\circ$，底部虚线上标 $2V\Delta t$，右侧标 $V\Delta t$、$\frac{\sqrt2}{2}V\Delta t$、角 $\alpha$，由 $M$ 向下的箭头 $V_M$，数条虚线(延长线)
% 图上方的文字“小量分析，延长 找角。”在图右上方(已在上面转录)

\noindent 用余弦定理\quad 勾股定\unclear{理}% 该行位于页面右上，行末被照片边缘截断

\[
V_M=\frac{\sqrt{\left(\sqrt2V\Delta t\right)^{2}+\left(\dfrac{\sqrt2}{2}V\Delta t\right)^{2}}}{\Delta t}=\frac{\sqrt{10}}{2}V
\]% 式末“V”右下有一小圆圈，像句号“。”，也可能是下标0(V_0)，未转录
%%%END
%%%BLOCK id=042-03 ch=04 sec=关联速度·伸缩等速 kind=method alt=none cont=none y=0.27-0.38
%FIG p042_e 0.02 0.277 0.30 0.375
\notefig[0.3]{figures/p042_e.png}
% 图内标注：左端小圆圈处标 $V_1$(向左箭头及虚线分量箭头)，细绳自左端圆圈斜向下到中间点、再斜向上连到右端小圆圈；右段绳上标 $V_3$(沿绳虚线箭头)；中间点处标 $V$(水平向右箭头)、$V_2$(斜向下箭头)
\begin{keybox}[伸缩等速]% 红笔框
\[
V_1+V_2=V_3\qquad \sum V_{\text{伸长}}=\sum V_{\text{缩短}}
\]
\end{keybox}
%%%END
%%%BLOCK id=042-04 ch=04 sec=关联速度 kind=method alt=03 cont=none y=0.25-0.38
\red{？}% 红笔问号，位于“两端圆周速度相等”左侧

\noindent 两端圆周速度相等\quad 一绳上 $T$、$x$、$v$、$a$、\unclear{$l$} 全相等% “a”后的一字像“1/I/l”

\noindent 绳不转动要平动

%FIG p042_f 0.668 0.297 0.835 0.376
\notefig[0.25]{figures/p042_f.png}
% 图内标注：斜板(与水平夹角 $30^\circ$)，板上标 $V$(沿板向右上箭头)，板与竖直绳交点处标“板 $\frac12V$”(向上箭头)、$\frac{\sqrt3}{2}V$(水平向右箭头)；竖直绳下端方块，绳上标 $V_1$(向上箭头)，方块处水平向右箭头
\noindent $\dfrac{\sqrt3}{2}V$ 抵消\unclear{车}转动% 位于图下方右侧(第三字像“车”，疑为“轮”)

\begin{align*}
V+\tfrac12V&=V_1\\
V_1&=\tfrac32V\\
\therefore V_{\text{合}}&=\sqrt3\,V
\end{align*}
%%%END
%%%BLOCK id=042-05 ch=04 sec=关联速度·弹性绳 kind=method alt=none cont=none y=0.38-0.50
\textbf{弹性绳}\ —\ 弹簧/激光/光线/轻杆套环。

%FIG p042_g 0.075 0.398 0.305 0.49
\notefig[0.28]{figures/p042_g.png}
% 图内标注：水平横线与竖直线、斜线(由竖线下端斜向上交于横线)；交点处黑字 $V_{\text{伸长}}$、$V_{\text{圆周}}$、$\omega$(弯箭头)；红字 $V\sin\theta$、$V$(向右箭头)、$V\cos\theta$，另有一小红字(像“沿”)

\noindent \underline{弹性物质无速度关联}。% 红笔下划线，并延伸为红箭头指向右侧红框

\begin{keybox}% 红笔框
%FIG p042_h 0.612 0.389 0.745 0.462
\notefig[0.25]{figures/p042_h.png}
% 图内标注：竖直线与水平线、弹簧(锯齿线)斜连在右上方小圆环与左下方小球之间(弹簧长度与角度都在变)
长度角度都变化、无法分解。
\end{keybox}
%%%END
%%%BLOCK id=042-06 ch=04 sec=关联速度·双绳拉船 kind=method alt=none cont=none y=0.50-0.66
\textbf{双绳拉船。}

%FIG p042_i 0.03 0.53 0.195 0.645
\notefig[0.2]{figures/p042_i.png}
% 图内标注：两只船 $A$、$B$(各带向外上方箭头)，用绳呈“Y”形系于下方物体 $C$，$C$ 上标 $V_C$(向上箭头)
%FIG p042_j 0.195 0.528 0.395 0.662
\notefig[0.28]{figures/p042_j.png}
% 图内标注：圆上三点 $A$(左)、$B$(右)、$C$(下)，连 $AB$、$AC$、$BC$，$C$ 处红字角 $\alpha$，竖直线由 $C$ 向上到圆顶并标 $V_C$(箭头)
\noindent 连$AB$\ 过$\overset{C}{B}$作画圆\qquad $V_C=\dfrac{|AB|}{\sin\alpha}$% “B”上方另有一小字母“C”(旁注/插入)，“B”笔画有叠写

\[
AB^{2}=\sqrt{V_A^{2}+V_B^{2}-2V_AV_B\cos\alpha}
\]% CHECK: 原稿等号左边写作 AB^2(上标2)，右边带根号，应为 AB=\sqrt{…} 或 AB^2 无根号，照原样转录
%%%END
%%%BLOCK id=042-07 ch=04 sec=关联速度·接触关联 kind=knowledge alt=none cont=none y=0.57-0.70
\textbf{接触关联}（左侧花括号分三条）
\begin{itemize}
\item 点点接触\quad \underline{不分离 $\Leftrightarrow$ $V$大小方向相同}
\item 点面接触\quad \underline{不分离 $\Leftrightarrow$ 垂直于接触面}\unclear{…}% 行末被照片边缘截断，其后可能还有字；箭头写作“(=)”形(手写 ⇔)，按 $\Leftrightarrow$ 转录
\item 面面接触\quad \underline{不分离} $\Leftrightarrow$ \underline{方向速度相同}
\end{itemize}
%%%END
%%%BLOCK id=042-08 ch=04 sec=关联速度·接触关联 kind=derivation alt=none cont=none y=0.69-0.80
%FIG p042_k 0.01 0.685 0.25 0.78
\notefig[0.25]{figures/p042_k.png}
% 图内标注：斜杆(下端小圆圈为转轴 $O$，与水平夹角 $\theta$，杆长 $L$)，杆上端 $A$，杆上沿杆箭头 $V_{B\parallel}$；杆与方块上角接触点 $B$，$V_B$(水平向右)、垂直杆的分速度 $V_{B\perp}$(斜向下箭头)、$M$(方块下沿处)、方块高 $h$
\begin{align*}
V_M&=V_B\\
V_{B\perp}&=V_B\sin\theta=\omega\cdot|OB|\\
\therefore\omega&=\frac{V\sin\theta}{|OB|}=\frac{V\sin\theta}{\dfrac{h}{\sin\theta}}=\frac{V\sin^{2}\theta}{h}\\
V_A&=L\cdot\omega=\frac{LV\sin^{2}\theta}{h}
\end{align*}% 原稿 ω 写得像 W
%%%END
%%%BLOCK id=042-09 ch=04 sec=关联速度·滑轮组位移关联 kind=problem alt=06 cont=none y=0.82-0.99
\red{？}% 红笔大问号，位于该题图左上方

%FIG p042_l 0.10 0.82 0.41 0.985
\notefig[0.3]{figures/p042_l.png}
% 图内标注：竖直墙面(左)、水平细线与两个小滑轮(圆圈)，水平段向左箭头；斜面(倾角 $\alpha$)上有物块 $B$(标 $m$，沿斜面向下的箭头)，红笔：$B$ 红字、红线连接滑轮与物块、红字 $X_B$ 及红色三角；斜面下方 $m$、$A$、$\alpha$，虚线；图下方“光滑”
\noindent $A$ 位移为 $x$ 时，求 $S_B$、$V_A$

\noindent 对$B$ 先向左，再沿斜面向下。

%FIG p042_m 0.435 0.918 0.577 0.97
\notefig[0.2]{figures/p042_m.png}
% 图内标注：位移矢量图——水平向左的 $x$(虚线)、沿斜面向下的 $x$、两段夹角红字 $\alpha/2$ 等，红粗线为合位移(斜向下箭头)
\[
S_B=2x\sin\frac{\alpha}{2}
\]

\[
mgx\sin\alpha=\frac12mV_x^{2}+\frac12m\left(V_x\cdot2\sin\frac{\alpha}{2}\right)^{2}
\]

\[
\therefore V_x=\sqrt{\frac{2gx\sin\alpha}{3-2\cos\alpha}}
\]
%%%END
%%%BLOCK id=043-01 ch=04 sec=关联速度·几何关联 kind=knowledge alt=none cont=none y=0.02-0.07
\begin{keybox}[几何关联]% 红笔框(页面左上角)，框内两行
$V$ 是三角形直角边
\end{keybox}
%%%END
%%%BLOCK id=043-02 ch=04 sec=关联速度·几何关联 kind=note alt=02,03 cont=none y=0.02-0.06
\red{受力要整体法}% 红色大字，位于页面右上角
%%%END
%%%BLOCK id=043-03 ch=04 sec=关联速度·几何关联 kind=problem alt=02,06 cont=none y=0.02-0.19
\begin{problem}[{[1]}]
%FIG p043_a 0.04 0.015 0.29 0.13
\notefig[0.3]{figures/p043_a.png}
% 图内标注：图左上角红框“几何关联 / V是三角形直角边”(已在上一块转录)；楔形斜面(左下角 $\theta$)，斜面上顶着一根竖直杆，杆顶标 $M$(重物)；楔形右侧标 $m$，水平向左箭头 $F$；红字 $V_M$(竖直向下，右侧)、$V_m$(水平，楔块下方红线箭头)
① 为将重物顶起，求 $F_{\min}$

\noindent $F_N=\dfrac{Mg}{\cos\theta}$\qquad $F_N\sin\theta=\dfrac{Mg}{\cos\theta}\sin\theta$

\noindent 将小轮推高地面 $h$ 处，突然撤去 $F$，求当小轮到水平面时斜面速度

\noindent ②
\[
\left\{\begin{aligned}
&V_M=V_m\tan\theta\\
&Mgh=\tfrac12MV_M^{2}+\tfrac12mV_m^{2}
\end{aligned}\right.
\]
\end{problem}
%%%END
%%%BLOCK id=043-04 ch=04 sec=关联速度·几何关联 kind=problem alt=02,06 cont=none y=0.19-0.44
\begin{problem}[{[2]}]
%FIG p043_b 0.09 0.19 0.345 0.325
\notefig[0.3]{figures/p043_b.png}
% 图内标注：竖直墙(左，顶端标 1)与斜面体(倾角 $30^\circ$，斜面上方标 2、$u=0$)之间夹一个小球(半径 $R$，标 $m$、两处 $R$)，球心到切点的虚线；斜面体右侧竖直边标 $M$，向左箭头 $F$，右侧红色小箭头 $V_1$，底边红箭头 $V_2$，右下标 $u=0$
撤去 $F$ 后小球和斜面体做匀加速直线运动，直到小球落地。

\noindent ① 求 $F$\quad ② 球落地时求 $V_1$、$V_2$

\noindent $F=N=\dfrac{\sqrt3}{3}mg$

\noindent $mgh=\tfrac12mV_1^{2}+\tfrac12mV_2^{2}$

\noindent $h_0=\sqrt3R$\quad $h=h_0-R=(\sqrt3-1)R$

\noindent $x_1=x_2\tan\theta$

\noindent $V_2=\sqrt3V_1$

\noindent $V_1=\sqrt{\dfrac{(\sqrt3-1)gR}{2}}$\quad $V_2=\sqrt{\dfrac{3(\sqrt3-1)gR}{2}}$
\end{problem}
%%%END
%%%BLOCK id=043-05 ch=04 sec=关联速度·几何关联 kind=problem alt=02,06 cont=none y=0.44-0.78
\begin{problem}[{[3]}]
%FIG p043_c 0.065 0.447 0.335 0.62
\notefig[0.3]{figures/p043_c.png}
% 图内标注：左侧楔形物块 $A$(顶角处标 $m$，下角 $37^\circ$、$\alpha$，红字 $V_B$ 水平向右箭头在顶端，红字 $V_A$ 竖直向下箭头在左侧)，紧靠的方块 $B$(标 $2m$)，弹簧(标 $k$)连接 $B$ 与右侧竖直墙；图下标 $\mu=0$；图左上角有题号 [3]
$B$ 将弹簧压缩 $x_0$ 后 $A$ 上滑，弹簧恢复原长时，$B$ 速度为 $V_B$。

\noindent ① 求刚恢复原长时 $V_A$。

\noindent $V_B=V_A\tan\alpha$\qquad $V_A=\dfrac43V_B$

%FIG p043_d 0.635 0.50 0.715 0.562
\notefig[0.25]{figures/p043_d.png}
% 图内标注：矢量图——向上箭头 $V_A$，向左箭头 $V_B$，夹角 $37^\circ$，虚线(斜向左上)
\noindent ② 求弹簧压缩量为 $x_0$ 时具有的弹性势能

\noindent 弹簧原长时，$h_A=\dfrac{x_0}{\tan\alpha}=\dfrac43x_0$

\noindent $E_p=\tfrac12mV_A^{2}+\dfrac{2}{2}mV_B^{2}+mgh=\dfrac{17}{9}mV_B^{2}+\dfrac43mgx_0$

\noindent ③ 求两物体动能最大时，弹簧压缩量

\noindent 设$AB$相互作用力$F$
\[
\left\{\begin{aligned}
&a_A=a_B=0\\
&F\sin\alpha=mg\\
&F\cos\alpha=kx
\end{aligned}\right.\qquad x=\frac{mg}{k\tan\alpha}=\frac{4mg}{3k}
\]
\end{problem}
%%%END
%%%BLOCK id=044-01 ch=04 sec=平抛运动·基本公式 kind=formula alt=none cont=none y=0.02-0.12
%FIG p044_a 0.262 0.022 0.535 0.115
\notefig[0.3]{figures/p044_a.png}
% 图内标注：斜面(下端标 $\theta$、$\alpha$)，抛体的抛物线轨迹；上端抛出点处 $V_0$(向左箭头)、$V\sin\theta$(向上箭头)，沿斜面 $V\cos\theta$；加速度分解 $g\sin\theta$、$g$(向下箭头)、$g\cos\theta$(斜向下箭头)，轨迹中部有小球及 $\theta$ 标记
\begin{align*}
x&=V_0t & V_x&=V_0\\
y&=\tfrac12gt^2 & V_y&=gt\\
S&=\sqrt{x^2+y^2} & V&=\sqrt{V_x^2+V_y^2}\\
\tan\alpha&=\frac{y}{x}=\frac{gt}{2V_0} & \tan\theta&=\frac{V_y}{V_x}=\frac{gt}{V_0}
\end{align*}
%%%END
%%%BLOCK id=044-02 ch=04 sec=平抛运动·基本公式 kind=note alt=none cont=none y=0.02-0.12
%FIG p044_b 0.64 0.028 0.78 0.08
\notefig[0.25]{figures/p044_b.png}
% 图内标注：“类上抛”“匀加速”两行字，左侧一个向左上箭头，另有一条斜向左下箭头由“匀加速”指向左下(指向左侧斜面图)
%FIG p044_c 0.80 0.025 0.955 0.101
\notefig[0.25]{figures/p044_c.png}
% 图内标注：L 形箭头——水平向右箭头旁写“匀速直线”，竖直向下箭头下写“自由落体”
\noindent $\nwarrow$ 类上抛\quad 匀加速$\swarrow$% 二者在页面右上方，箭头指向左侧斜面图

\noindent $\rightarrow$ 匀速直线\quad $\downarrow$ 自由落体

\noindent \red{抛体\unclear{重}分解}

\noindent \red{平抛\unclear{有}两个直角三角形\ 一个反向延长线}

\noindent \red{斜抛\unclear{有}从顶部的两个平抛}% 以上三行均为红笔；“有”字辨认不清(像“考”)
%%%END
%%%BLOCK id=044-03 ch=04 sec=平抛运动·斜面平抛 kind=formula alt=none cont=none y=0.12-0.20
\textbf{① 落点时间}
\[
\tan\theta=\frac{gt}{2V_0}=\frac12\tan\alpha\qquad
\tan\theta=\frac{\frac12gt^2}{V_0t}\qquad
t=\frac{2V_0\tan\theta}{g}
\]
%%%END
%%%BLOCK id=044-04 ch=04 sec=平抛运动·重力功率 kind=formula alt=06 cont=none y=0.19-0.26
\textbf{② 重力功率}
\[
\left\{\begin{aligned}
\bar P_G&=\frac{mgy}{t}=\frac{mg\cdot\frac12gt^2}{t}=\frac12mg^2t\\
P_G&=mgV\cos\alpha=mgV_y=mg^2t
\end{aligned}\right.
\qquad
\left\{\begin{aligned}
\bar V&=V_{\frac t2}=\frac{V_0+V}{2}\\
\bar P&=P_{\frac t2}=\frac{P_0+P}{2}
\end{aligned}\right.
\]% 一条细弧线自 mgVcosα 附近向左下延伸至 ③ 的 h 式附近(指示线)
%%%END
%%%BLOCK id=044-05 ch=04 sec=平抛运动·斜面平抛 kind=formula alt=none cont=none y=0.25-0.33
\textbf{③ 离斜面最远}
\[
h=\frac{(V_0\sin\theta)^2}{2g\cos\theta}\qquad
t=\frac{V_0\tan\theta}{g}\qquad
V_{\parallel}=V_0\cos\theta+g\sin\theta\,t
\]
\noindent 半程$\nearrow$% 小字“半程”带箭头，指向 t 式
%%%END
%%%BLOCK id=044-06 ch=04 sec=平抛运动·斜面平抛 kind=method alt=01 cont=none y=0.32-0.41
\textbf{④}

%FIG p044_d 0.095 0.335 0.225 0.39
\notefig[0.25]{figures/p044_d.png}
% 图内标注：斜面上方的抛物线轨迹，沿斜面方向分成两段，分别标 1、2，右端标 $S_1$(带向左箭头)
\noindent $\dfrac{S_1}{S_2}=\dfrac13$\qquad $V_{y_0}=0$

%FIG p044_e 0.51 0.335 0.66 0.395
\notefig[0.25]{figures/p044_e.png}
% 图内标注：$V$-$t$ 图像(纵轴 $V$，横轴 $t$)，过原点的直线下三角形按相等时间段分成面积 1、3、5 三块
\noindent $V_{y_0}\neq0$ 时\quad 若为

%FIG p044_f 0.685 0.34 0.83 0.395
\notefig[0.25]{figures/p044_f.png}
% 图内标注：$V$-$t$ 图像(上部为过原点直线下的三角形，分成 1、3 两块；下部为两个相等矩形，各标 $S_0$)
\noindent \underline{真分数不等式。}% 红线之外为黑笔下划线

\[
\frac{1+S_0}{3+S_0}>\frac13
\]
\[
\therefore 1>\frac{S_1}{S_2}>\frac13
\]% 末式下方有波浪线
%%%END
%%%BLOCK id=044-07 ch=04 sec=平抛运动·轨迹方程 kind=formula alt=none cont=none y=0.41-0.50
\textbf{⑤ 轨迹方程}\quad $y=\dfrac{g}{2V_0^2}x^2$\quad (用于联立，建系) 找交点。
\[
\left\{\begin{aligned}x&=V_0t\\y&=\tfrac12gt^2\end{aligned}\right.
\]
%%%END
%%%BLOCK id=044-08 ch=04 sec=平抛运动·偏角 kind=formula alt=none cont=none y=0.50-0.60
\textbf{⑥ 偏角}
\[
\left\{\begin{aligned}
\tan\theta&=\frac{gt}{V_0}\quad\text{速偏}\\
\tan\alpha&=\frac{gt}{2V_0}\quad\text{位偏}
\end{aligned}\right.
\]

%FIG p044_g 0.46 0.482 0.61 0.545
\notefig[0.25]{figures/p044_g.png}
% 图内标注：水平线上两个反向箭头，左边标 $V_1$、右边标 $V_2$，中间竖直虚线；箭头旁标 1(向左)、2(向右)
\noindent 当 1、2 速度垂直时求 $t$
\[
\tan\theta_1\cdot\tan\theta_2=1
\]
\[
\tan\theta_1=\frac{gt}{V_1}\qquad\tan\theta_2=\frac{gt}{V_2}\qquad t=\frac{\sqrt{V_1V_2}}{g}
\]
%%%END
%%%BLOCK id=044-09 ch=04 sec=平抛运动·抛上另一斜面 kind=method alt=06 cont=none y=0.595-0.84
\textbf{⑦ 抛上另一个斜面}

%FIG p044_h 0.125 0.617 0.255 0.69
\notefig[0.25]{figures/p044_h.png}
% 图内标注：$V_0$(水平向右箭头)，抛物线落到斜面上，斜面倾角 $\theta$，落点处速度箭头及夹角 $\varphi$
\noindent 若垂直，$\tan\varphi=\dfrac{gt}{V_0}=\dfrac{1}{\tan\theta}$

\[
\bar P=\frac{mgy}{t}=\frac{mg\cdot\frac12gt^2}{t}=\frac{mgV_0\cot\theta}{2}=mgV_y
\]% CHECK: 末项 mgV_y 与前一项 mgV_0cotθ/2 相差系数 1/2，疑漏写或 V_y 指平均值，照原样转录；“P̄=”之后有一小块被涂白

%FIG p044_i 0.11 0.69 0.275 0.745
\notefig[0.25]{figures/p044_i.png}
% 图内标注：$V_0$(水平向右箭头)，斜面倾角 $\theta$，轨迹与斜面“最短距离”处(标“最短距离”并带引线)，速度与斜面夹角 $90^\circ-\theta$
\[
\frac{gt}{2V_0}=\tan(90^\circ-\theta)=\frac{1}{\tan\theta}
\]

%FIG p044_k 0.065 0.752 0.265 0.83
\notefig[0.25]{figures/p044_k.png}
% 图内标注：左侧花括号标 $H$(总高度)；$V_0$ 水平向右箭头，抛物线落到倾角 $\theta$ 的斜面上，落点速度方向 $\theta$，虚线辅助线；两条细长弧线由图中两段高度分别指向右侧式中 $V_0t\tan\theta$ 与 $\frac12gt^2$ 的下划线
\[
\left\{\begin{aligned}
H&=\underline{V_0t\tan\theta}+\underline{\tfrac12gt^2}\\
\frac{1}{\tan\theta}&=\frac{gt}{V_0}
\end{aligned}\right.
\qquad V_0=\sqrt{gH}
\]% CHECK: 由左式联立得 H=(V_0^2/g)(1+cot^2θ/2)，与 V_0=√(gH) 不一致，照原样转录
%%%END
%%%BLOCK id=044-10 ch=04 sec=平抛运动·偏角之差(不单调) kind=method alt=none cont=none y=0.70-0.90
\textbf{⑧}

%FIG p044_j 0.605 0.70 0.75 0.78
\notefig[0.25]{figures/p044_j.png}
% 图内标注：斜面(倾角 $\varphi$)，上方箭头轨迹，标 $\alpha$、$\theta$、$\varphi$ 等角，水平虚线
\[
\tan\theta=\tan(\alpha-\varphi)=\frac{\tan\alpha-\tan\varphi}{1+\tan\alpha\tan\varphi}
\]
\[
\tan\alpha=2\tan\varphi
\]
\[
\therefore\tan\theta=\frac{1}{2\tan\varphi+\dfrac{1}{\tan\varphi}}
\]

\noindent 不单调
%FIG p044_l 0.78 0.803 0.965 0.895
\notefig[0.25]{figures/p044_l.png}
% 图内标注：坐标轴(纵轴带箭头)与一条过原点的波形曲线：原点左侧先下凹，右侧上凸后下降，旁写“不单调”
%%%END
%%%BLOCK id=044-11 ch=04 sec=平抛运动·无碰撞 kind=method alt=none cont=none y=0.84-0.99
\noindent 无碰撞 = 速偏 $\Leftrightarrow$ $V\parallel$接触面% “接触面”三字写得较连，辨认为“接触面”

%FIG p044_m 0.09 0.872 0.225 0.95
\notefig[0.25]{figures/p044_m.png}
% 图内标注：$\Rightarrow$ 水平箭头，虚线轨迹，到斜面上的小物体(管口)，弯箭头；图下标“平行”
%FIG p044_n 0.225 0.87 0.385 0.955
\notefig[0.25]{figures/p044_n.png}
% 图内标注：“0$\Rightarrow$”水平箭头，虚线轨迹，进入半圆形碗(右上是碗口)，向下箭头及虚线；图下标“切向”
\noindent 平行\qquad 切向

%FIG p044_o 0.435 0.864 0.605 0.95
\notefig[0.25]{figures/p044_o.png}
% 图内标注：$\Rightarrow V_0$ 水平箭头，虚线轨迹落到半圆弧上，圆弧上标 $60^\circ$、$30^\circ$、$120^\circ$(像“160°”)，末端速度箭头 $V$
\noindent 刚好相切无碰

\[
\left\{\begin{aligned}
\tan30^\circ&=\frac{gt}{V_0}\\
\frac{\sqrt3}{2}R&=V_0t
\end{aligned}\right.
\]% CHECK: “R”前分数的分子写得像“3”(根号不明)，按 √3/2 转录

\noindent \struck{$\therefore V_0=3\sqrt5\unit{m/s}$}% 整行被横线划去

\noindent $\Rightarrow V_0$
%%%END
%%%BLOCK id=045-01 ch=04 sec=平抛运动·斜面平抛 kind=knowledge alt=none cont=none y=0.015-0.10
%FIG p045_a 0.08 0.015 0.185 0.095
\notefig[0.25]{figures/p045_a.png}
% 图内标注：斜面，抛出点处水平箭头 $V_0$，几条抛物线落到斜面上，标 $2\alpha$、$\alpha$ 等角
\noindent 位偏 = 从斜到斜 = 位偏角为倾斜角\qquad $\tan\alpha=\dfrac{gt}{2V_0}$
%%%END
%%%BLOCK id=045-02 ch=04 sec=平抛运动·斜面平抛 kind=method alt=none cont=none y=0.115-0.42
%FIG p045_b 0.095 0.115 0.31 0.23
\notefig[0.3]{figures/p045_b.png}
% 图内标注：同一抛出点，两个水平初速度 $V_0$、$2V_0$ 的抛物线落到斜面上，红色直线连抛出点与两个落点，落点处标 $\alpha_1$、$\alpha_2$，右下标 $S_1$ 并写“求 $S_1/S_2$”
\noindent 求 $\dfrac{S_1}{S_2}$

\noindent 讨论\quad ① 都落在斜面上

\noindent $\tan\alpha=\dfrac{gt}{2V_0}$\qquad $t=\dfrac{2V_0\tan\alpha}{g}\propto V_0$\qquad \fbox{$t_1:t_2=1:2$}

\noindent $S=V_0t\propto V_0^2$\qquad \fbox{$\dfrac{S_1}{S_2}$ 则为 $1:4$}% 方框内“则”字辨认略不确定

\noindent ② 都落在平面

\noindent $h=\tfrac12gt^2$\quad \fbox{$t_1=t_2$}\quad $S=V_0t\propto V_0$\quad \fbox{$\dfrac{S_1}{S_2}=1:2$}

\noindent ③ 一个斜面一个平面

\noindent $\alpha_1>\alpha_2$

\noindent $\therefore\tan\alpha_1>\tan\alpha_2$

\noindent $\therefore\dfrac{gt_1}{2V_0}>\dfrac{gt_2}{2V_0}$

\noindent \fbox{$\begin{array}{l}\therefore\dfrac12<\dfrac{t_1}{t_2}<1\\[4pt]\therefore\dfrac14<\dfrac{S_1}{S_2}<\dfrac12\end{array}$}% 该方框位于 ③ 推导式右侧
%%%END
%%%BLOCK id=045-03 ch=04 sec=平抛运动·斜面平抛 kind=method alt=none cont=none y=0.42-0.545
%FIG p045_c 0.055 0.42 0.30 0.492
\notefig[0.3]{figures/p045_c.png}
% 图内标注：两个对称斜面构成 V 形，各高 $h$、水平长 $2h$，左斜面顶端抛出 $V$(水平箭头)，抛物线落到谷底 $O$ 点；$O$ 处标 $\theta_1$、$\theta_2$，向下箭头
\noindent $t_{\text{最短}}$\quad $t=\sqrt{\dfrac{2h}{g}}$

\noindent 找界点：在 $O$ 点\quad $\theta_1<45^\circ$\qquad $\tan\theta_2=2\tan\theta_1=\dfrac{2y}{x}=1$

\noindent $\theta_1+\theta_2<90^\circ$\qquad 故无法垂直打在斜面 (若末速度反向\unclear{延}长线过位移中点。)% 括号内文字位于页面右侧，第一行末“过”字靠近页边，第二行为“位移中点。)”
%%%END
%%%BLOCK id=045-04 ch=04 sec=平抛运动·斜面平抛 kind=problem alt=06 cont=none y=0.545-0.70
%FIG p045_d 0.09 0.553 0.355 0.695
\notefig[0.3]{figures/p045_d.png}
% 图内标注：左上小圆圈出发水平箭头 $V_0$(抛出点)；斜面(高 $h$，带花括号)，斜面上沿斜面向下箭头 $V_1$；斜面底端连水平面，水平面上标 $x$、$\mu$；图下斜写小字“如履平地”(第二字像“覆”)
\noindent 如\unclear{履}平地% 斜写于图下方

\noindent 求能走多远
\[
\left\{\begin{aligned}
V_1&=V_0\cos\theta+gt\sin\theta\\
\tan\theta&=\frac{gt}{2V_0}\\
h&=H-\tfrac12gt^2\\
\tfrac12mV_1^2&=mgh-\mu mg\cos\theta\,\frac{h}{\sin\theta}-\mu mgx
\end{aligned}\right.
\]
%%%END
%%%BLOCK id=045-05 ch=04 sec=圆周运动·竖直面临界 kind=summary alt=06 cont=none y=0.71-0.85
%FIG p045_e 0.055 0.708 0.21 0.84
\notefig[0.25]{figures/p045_e.png}
% 图内标注：每行左侧的示意小图——第一行红笔：竖线上端带小圆(绳)、半圆弧∪(槽)、小方框(筒)；第二行：红色竖线带小圈(杆)，“环”“管”二字各被红圈圈出；第三行：黑色拱形∩
\noindent $mg+T=\dfrac{mV^2}{r}$% 位于表格上方、第一行方程栏

\noindent \red{$V<\sqrt{gr}$ 或 $\omega<\sqrt{\dfrac{g}{r}}$ 则对筒槽绳压力为0 即为最小值}% 红笔，位于第一行上方

\begin{center}
\small
\begin{tabular}{@{}l l l c c c c@{}}
\toprule
类型 & 受力特点 & 最高点方程 & 最高 & 最低 & 上下两点压力差$\Delta N$ & 向心力上下差\\
\midrule
绳\ 槽\ 筒 & 无法提供支持力、给压力 & $mg+T=\dfrac{mV^2}{r}$ & $V=\sqrt{gR}$ & $V=\sqrt{5gR}$ & $6mg$ & $4mg$\\[4pt]
杆\ 环\ 管 & 双向力 & $mg\pm F_N=\dfrac{mV^2}{r}$ & $0$ & $V=\sqrt{4gR}$ & $6mg$ & $4mg$\\[4pt]
拱 & 支持力 & $mg-N=\dfrac{mV^2}{r}$ & $0$ & 无 & 无 & 无\\
\bottomrule
\end{tabular}
\end{center}% 原稿表格左侧有一个大花括号括住三行；“类型/受力特点/最高点方程”三个表头为转录者补写的栏名(原稿无)；第一行的方程 mg+T=mV^2/r 在原稿中位于该行上方；第三行方程前有一小块被涂白
%%%END
%%%BLOCK id=045-06 ch=04 sec=圆周运动·竖直面临界 kind=method alt=06 cont=none y=0.85-0.99
\noindent ☆ 不脱离、\underline{绷直}

%FIG p045_f 0.285 0.85 0.395 0.945
\notefig[0.2]{figures/p045_f.png}
% 图内标注：圆(半径 $R$)，圆心到右侧等高点的水平虚线、到最低点的竖直虚线；圆上标 $V_2$(最高点)、$V_1$(与圆心等高的右侧点)；最低点有一个小圆圈(出发点)
\noindent $V\le V_1$ 或 $V\ge V_2$

\noindent $V\le\sqrt{2gR}$\qquad $V\ge\sqrt{5gR}$
%%%END
%%%BLOCK id=046-01 ch=04 sec=圆周运动·圆心参考系 kind=knowledge alt=none cont=prev y=0.03-0.13
$V$为相对圆心的速度，而非单单相对地面

圆心参考系．\ $F_n=\sum F_{\text{指向圆心}}-\sum F_{\text{背离圆心}}$\ （\unclear{分解}）

\[ F_n=\frac{mv^2}{R}=\frac{2E_k}{R} \]
%%%END
%%%BLOCK id=046-02 ch=04 sec=圆周运动·传动 kind=knowledge alt=none cont=none y=0.08-0.14
传动\ $\begin{cases}\text{同点同}\,V\\ \text{同轴同}\,\omega\end{cases}$
%%%END
%%%BLOCK id=046-03 ch=04 sec=圆周运动·竖直面临界 kind=method alt=06 cont=none y=0.14-0.44
\textbf{技巧}

%FIG p046_a 0.08 0.14 0.30 0.305
\notefig[0.3]{figures/p046_a.png}

\[
\begin{cases}
mg\,2R+W_{\text{其他}}=\dfrac{1}{2}mV_{\text{下}}^{2}-\dfrac{1}{2}mV_{\text{上}}^{2}\\[4pt]
N_{\text{上}}+mg=\dfrac{mV_{\text{上}}^{2}}{R}\\[4pt]
N_{\text{下}}-mg=\dfrac{mV_{\text{下}}^{2}}{R}
\end{cases}
\]
\[ \therefore\ N_{\text{下}}-N_{\text{上}}=6mg+\frac{2W_{\text{其他}}}{R} \]
\[ \Delta N=6mg+\frac{2W_{\text{其他}}}{R} \]
\[ \text{实际}\,\Delta N=\text{理想}\,\Delta N+2\,\frac{W_{\text{其他}}}{R} \]
\[ \text{偏差}\,\Delta N\,\frac{R}{2}=W_{\text{其他}}\qquad W_{\text{其他}}=\frac{\Delta N_{\text{偏差}}}{2}R\qquad \Delta N_{\text{偏差}} \]
%%%END
%%%BLOCK id=046-04 ch=04 sec=运动合成·相对速度 kind=knowledge alt=none cont=none y=0.44-0.60
\textbf{运动合成}\ \ {\footnotesize $P$，$I$，$V$，$a$，$X$}

\[ ab^{\text{对}}=a_{\text{对地}}-b_{\text{对地}} \]
% CHECK: 公式左端"αb"及其右上角的"对"辨认不清；右端"对地"的"对"写在"地"字上方；"运动合成"右上方有小字 P、I、V、a、X（疑为动量、冲量、速度、加速度、位移皆为相对量）

%FIG p046_b 0.115 0.505 0.435 0.605
\notefig[0.5]{figures/p046_b.png}
% 图 p046_b 内含(黑字)公式 V雨车=V雨地−V车地 及矢量三角形、各标注

\[ V_{\text{雨车}}=V_{\text{雨地}}-V_{\text{车地}} \]

%FIG p046_d 0.50 0.495 0.685 0.59
\notefig[0.3]{figures/p046_d.png}

\[ V_{\text{刀玻}}=V_{\text{刀地}}-V_{\text{玻地}} \]
%%%END
%%%BLOCK id=046-05 ch=04 sec=小船过河 kind=method alt=none cont=none y=0.60-0.97
\textbf{小船过河}\quad ① 最快\ $t=\dfrac{d}{V_{\text{船}}}$\quad ② 最近\ $V_{\text{船}}>V_{\text{水}}$\ 垂直过河

$V_{\text{船}}<V_{\text{水}}$\quad $V_{\text{水}}$为斜边，$V_{\text{船}}$画圆找切线\quad $\sin\alpha=\dfrac{d}{L}=\dfrac{V_{\text{船}}}{V_{\text{水}}}$

%FIG p046_e 0.70 0.641 0.97 0.70
\notefig[0.45]{figures/p046_e.png}

%FIG p046_f 0.285 0.675 0.445 0.76
\notefig[0.3]{figures/p046_f.png}

甲乙$V$相等\unclear{向}P，在NP上\unclear{相}遇，渡河时间一致

%FIG p046_g 0.445 0.733 0.56 0.795
\notefig[0.25]{figures/p046_g.png}

$V_2$相对甲

%FIG p046_h 0.05 0.82 0.297 0.915
\notefig[0.35]{figures/p046_h.png}

船头去\unclear{时}垂直河岸，归\unclear{时}路线垂直河岸，

\[ t_1=\frac{d}{V_{\text{船}}}\qquad t_2=\frac{d}{\sqrt{V_{\text{船}}^{2}-V_{\text{河}}^{2}}}\qquad \frac{t_1}{t_2}=k\qquad k=\sqrt{1-\left(\frac{V_{\text{水}}}{V_{\text{船}}}\right)^{2}} \]
% CHECK: t_2 分母根号内 V船² 与 V河² 之间的符号原稿像"="（疑为"−"的笔误），此处按"−"录入；t_2 中写 V河、k 中写 V水，照原稿
%%%END
%%%BLOCK id=068-01 ch=04 sec=曲线运动·条件与速度变化 kind=knowledge alt=none cont=none y=0.07-0.25
% 原稿：页面最上方一行（位于“Date. NO.”表头上方），以及表头下方右侧两行、其下两处批注
$a$\ $V$锐角\ $V\uparrow$\qquad $a$\ $V$钝角\ $V\downarrow$\qquad $a$\ $V$\unclear{垂直}\ 速度大小不变，方向改变 % CHECK: 顶行三项中 a 与 V 之间似无符号（应为夹角关系），第三项“垂直”二字不确定

$a\perp V$\quad 改变$V$方向\quad 法向加速度

$a\parallel V$\quad 改变$V$大小\quad 切向加速度

$a$、$V$不同方向

力与速度夹轨迹\quad 凹侧方向偏向力
%%%END
%%%BLOCK id=068-02 ch=04 sec=章节提纲 kind=summary alt=none cont=none y=0.12-0.97
% 版面：左侧一列为各部分标题（曲线运动；运动的合成与分解，其下一行小字“合成、分解”；抛体运动，其下“分解”；圆周运动，其下一行小字），右侧为对应条目
\begin{itemize}
  \item \textbf{曲线运动}
  \begin{itemize}
    \item 曲线运动、骑马射箭、转台投篮、风中骑车、\unclear{滑}板运动
  \end{itemize}
  \item \textbf{运动的合成与分解}\\ 合成、分解
  \begin{itemize}
    \item 运动的合成与分解，\quad 等时性\quad 独立性\quad 等效性
    \item 小船渡河模型、
    \item 关联速度模型、 % 图及批注见 068-03
  \end{itemize}
  \item \textbf{抛体运动}\\ 分解
  \begin{itemize}
    \item （类）平抛运动
    \item （类）斜抛运动
    \item 落点有约束的平抛运动
    \item 平抛运动的临界问题
    \item 斜抛运动问题。
  \end{itemize}
  \item \textbf{圆周运动}\\ \unclear{受力}
  \begin{itemize}
    \item 线速度\quad \underline{角速度}\quad \underline{周期}、\underline{转速}\quad 向心加速度\quad 向心力
    \item 离心运动（$F_n<m\omega^2 r$）、和向心运动（$F_n>m\omega^2 r$） % 原稿两式分别写在“离心运动”“和向心运动”上方
    \item $\underbrace{\text{转杆}}_{\text{同}\,\omega}$、$\underbrace{\text{皮带传动、齿轮传动、摩擦传动}}_{\text{同}\,v}$、$\underset{\text{同}\,\omega}{\text{同轴传动}}$ % 原稿批注写在各项下方：转杆下波浪线+“同ω”；皮带/齿轮/摩擦传动下长波浪线+“同v”；同轴传动下“同ω”
  \end{itemize}
\end{itemize}
%%%END
%%%BLOCK id=068-03 ch=04 sec=关联速度模型 kind=diagram alt=none cont=none y=0.34-0.46
\textbf{关联速度模型}

%FIG p068_a 0.335 0.346 0.494 0.415
\notefig[0.28]{figures/p068_a.png}

%FIG p068_b 0.49 0.366 0.654 0.45
\notefig[0.28]{figures/p068_b.png}

沿绳垂绳，沿杆垂杆

力和速度的分解不同
%%%END
%%%BLOCK id=069-01 ch=04 sec=水平面圆周运动·飞机转弯 kind=method alt=none cont=none y=0.09-0.255
\textbf{水平面内的匀速圆周运动。}

\textbf{飞机水平转弯}

%FIG p069_a 0.29 0.16 0.42 0.255
\notefig[0.25]{figures/p069_a.png}

\[ F_n=mg\tan\theta \]
%%%END
%%%BLOCK id=069-02 ch=04 sec=水平面圆周运动·火车转弯 kind=method alt=none cont=none y=0.25-0.345
\textbf{火车转弯}

%FIG p069_b 0.26 0.2545 0.404 0.345
\notefig[0.25]{figures/p069_b.png}

\[ F_n=mg\tan\theta \]
%%%END
%%%BLOCK id=069-03 ch=04 sec=水平面圆周运动·飞车走壁 kind=method alt=none cont=none y=0.36-0.45
\textbf{飞车走壁}

%FIG p069_c 0.22 0.36 0.47 0.45
\notefig[0.4]{figures/p069_c.png}

\[ F_n=mg\tan\theta \]
%%%END
%%%BLOCK id=069-04 ch=04 sec=水平面圆周运动·汽车转弯 kind=method alt=none cont=none y=0.455-0.51
\textbf{汽车转弯}

%FIG p069_d 0.19 0.455 0.31 0.51
\notefig[0.25]{figures/p069_d.png}

\[ f=F_n \]
%%%END
%%%BLOCK id=069-05 ch=04 sec=水平面圆周运动·圆筒内壁 kind=method alt=none cont=none y=0.455-0.57
%FIG p069_e 0.49 0.449 0.63 0.535
\notefig[0.25]{figures/p069_e.png}

\begin{align*}
& mg=f\le\mu F_N\\
& F_N=m\omega^2 r\\
& \omega\ge\sqrt{\dfrac{g}{\mu r}}
\end{align*}
转速越大越不易掉下
%%%END
%%%BLOCK id=069-06 ch=04 sec=水平面圆周运动·水平转台 kind=method alt=none cont=none y=0.535-0.62
\textbf{水平转台}

%FIG p069_f 0.20 0.535 0.35 0.62
\notefig[0.25]{figures/p069_f.png}

\[ F_n=m_Bg \]
%%%END
%%%BLOCK id=069-07 ch=04 sec=水平面圆周运动·圆锥摆 kind=method alt=none cont=none y=0.58-0.80
\textbf{圆锥摆}

%FIG p069_g 0.17 0.645 0.35 0.78
\notefig[0.3]{figures/p069_g.png}

\begin{align*}
& T\cos\theta=mg\\
& T\sin\theta=mg\tan\theta=F_{\text{向}}
\end{align*}

\textbf{单个圆锥摆}
\begin{align*}
F_n&=mg\tan\theta & v&=\sqrt{ar}=\sqrt{g\tan\theta\sin\theta}\quad\text{与}\theta\text{正相关}\\ % CHECK: 根号内疑似漏写 L（r=L sin θ 时 v=sqrt(gL tanθ sinθ)），照原样转录
r&=L\sin\theta & \omega&=\sqrt{\dfrac{a}{r}}=\sqrt{\dfrac{g}{H}}\\
T&=mL\omega^2=\dfrac{mg}{\cos\theta} & a_{\text{向}}&=g\tan\theta=\omega^2r=\dfrac{v^2}{r}
\end{align*}
%%%END
%%%BLOCK id=069-08 ch=04 sec=水平面圆周运动·圆锥摆 kind=derivation alt=none cont=none y=0.79-0.86
% 左侧页边小标题为“水平转盘、单列”（第二行两字不确定），其右为下列推导
\textbf{水平转盘、\unclear{单列}}
\[ mg\tan\theta=m\omega^2L\sin\theta=\dfrac{m\cdot4\pi^2}{T^2}L\sin\theta,\qquad \omega=\sqrt{\dfrac{g}{L\cos\theta}}=\sqrt{\dfrac{g}{H}} \]
\[ T=2\pi\sqrt{\dfrac{L\cos\theta}{g}}=2\pi\sqrt{\dfrac{H}{g}} \]
%%%END
%%%BLOCK id=069-09 ch=04 sec=水平面圆周运动·双圆锥摆 kind=method alt=none cont=none y=0.85-0.97
\textbf{双圆锥摆}

%FIG p069_h 0.30 0.854 0.49 0.985
\notefig[0.3]{figures/p069_h.png}

$\omega$相等，$V$和$a$、$\propto r$，

竖直方向合力为0
%%%END
%%%BLOCK id=070-01 ch=04 sec=竖直面圆周运动·轻绳模型 kind=method alt=none cont=none y=0.08-0.31
\textbf{竖直圆周运动}

\textbf{轻绳模型}

%FIG p070_a 0.14 0.148 0.355 0.192
\notefig[0.3]{figures/p070_a.png}

\begin{align*}
& mg+F_T=\dfrac{mv^2}{r}\\
& F_T=0\qquad mg=\dfrac{mv^2}{r}\\
& v=\sqrt{gr}\\
& F=\dfrac{mv^2}{r}-mg
\end{align*}

无支持力

绳力只能向下。

\[ V_{\text{过上}}=\sqrt{gr} \]
%%%END
%%%BLOCK id=070-02 ch=04 sec=竖直面圆周运动·轻杆模型 kind=method alt=none cont=none y=0.15-0.31
\textbf{轻杆模型}

%FIG p070_b 0.425 0.178 0.62 0.214
\notefig[0.3]{figures/p070_b.png}

弹力可向下\quad 可向上\quad 也为0

\unclear{当}\ $V=0$\quad $F_n=0$\quad $F_N=mg$

$V_{\text{过上}}=0$

%FIG p070_c 0.625 0.128 0.815 0.22
\notefig[0.3]{figures/p070_c.png}

\begin{align*}
& mg\pm F_N=\dfrac{mv^2}{r}\\
& F=\dfrac{mv^2}{r}-mg\\
& mg-(-F)=\dfrac{mv^2}{r}
\end{align*}
%%%END
%%%BLOCK id=076-01 ch=04 sec=水平面圆周·圆锥内壁 kind=method alt=none cont=none y=0.04-0.20
%FIG p076_a 0.10 0.045 0.225 0.17
\notefig[0.25]{figures/p076_a.png}
\[
\begin{aligned}
&N\sin\theta = mg\\
&N\cos\theta = mg\cot\theta\\
&a_{\text{向}} = g\cot\theta = \omega^2 r = \frac{v^2}{r}
\end{aligned}
\qquad\qquad
\begin{aligned}
&v=\sqrt{ar}=\sqrt{g\cot\theta\,r}\\
&\omega=\sqrt{\frac{a}{r}}=\sqrt{\frac{g}{H}}\\
&T=2\pi\sqrt{\frac{H}{g}}
\end{aligned}
\]
% CHECK: 图中 H 的含义不明（标在水平箭头旁），ω=√(g/H)、T=2π√(H/g) 与 a=g cotθ 的关系需核对
%%%END
%%%BLOCK id=076-02 ch=04 sec=水平面圆周·半球碗内壁 kind=method alt=none cont=none y=0.21-0.34
%FIG p076_b 0.07 0.215 0.29 0.34
\notefig[0.3]{figures/p076_b.png}
\[
\begin{aligned}
&N\cos\theta = mg\\
&N\sin\theta = mg\tan\theta = F_n\\
&r = R\sin\theta
\end{aligned}
\qquad\qquad
\begin{aligned}
&v=\sqrt{ar}=\sqrt{g\tan\theta\sin\theta}\quad\text{正相关}\\
&\omega=\sqrt{\frac{a}{r}}=\sqrt{\frac{g}{H}}=\sqrt{\frac{g\tan\theta}{L\sin\theta}}\\
&T=2\pi\sqrt{\frac{H}{g}}
\end{aligned}
\]
% CHECK: v=√(g tanθ sinθ) 手写如此，按 r=R sinθ 似漏写 R；ω 末式分母手写像 L，按 r=R sinθ 应为 R
%%%END
%%%BLOCK id=076-03 ch=04 sec=水平面圆周·半球碗内壁 kind=derivation alt=none cont=none y=0.37-0.60
%FIG p076_c 0.09 0.375 0.30 0.47
\notefig[0.3]{figures/p076_c.png}
\textbf{若光滑}
\begin{align*}
&N\cos\theta = mg\\
&N\sin\theta = mg\tan\theta = ma\\
&H = R(1-\cos\theta)\\
&a = g\tan\theta = \omega^2 r = \omega^2 R\sin\theta \qquad \cos\theta = \frac{g}{\omega^2 R}\\
&\omega\uparrow\quad H\uparrow\quad \omega\uparrow\ \text{上升．}
\end{align*}
% CHECK: 最后一行第三个符号手写同 ω（按推导可能是 θ↑）
%%%END
%%%BLOCK id=076-04 ch=04 sec=水平面圆周·半球碗内壁 kind=knowledge alt=none cont=none y=0.63-0.73
\textbf{若不光滑．}
\begin{center}
\begin{tabular}{@{}ll@{}}
$a=g\tan\theta$ & 不受 $f$\\
$a>g\tan\theta$ & $f$ 向下\\
$a<g\tan\theta$ & $f$ 向上\\
\end{tabular}
\end{center}
%%%END
%%%BLOCK id=087-08 ch=04 sec=竖直面圆周·系统牛二 kind=problem alt=03,06 cont=none y=0.72-0.90
\begin{problem}
%FIG p087_g 0.105 0.73 0.2 0.795
\notefig[0.25]{figures/p087_g.png}
在最高点恰好对地面无压力，求在最低点压力
\begin{solution}
\[ (M+m)g=ma_{\text{上}}=m\,\frac{V_{\text{上}}^2}{R} \]
\[ N-(M+m)g=\frac{mV_{\text{下}}^2}{R}\qquad \text{在圆周必用动能定理} \]
\[ \frac12mV_{\text{下}}^2-\frac12mV_{\text{上}}^2=mg\,2R \]
\[ \therefore N=(2M+6m)g \]
\end{solution}
\end{problem}
%%%END
%%%BLOCK id=088-03 ch=04 sec=绳拉直问题 kind=method alt=06 cont=none y=0.67-0.95
\noindent \textbf{绳拉直问题}\quad 沿绳方向速度消失
% “绳拉直问题”原稿被椭圆圈住
%FIG p088_d 0.09 0.735 0.41 0.93
\notefig[0.45]{figures/p088_d.png}
\[ \begin{cases} V^2=2gL\sin\theta\\ mg(L-L\sin\theta)=\dfrac12mV'^2-\dfrac12m(V\cos\theta)^2\end{cases} \]
% CHECK: 第一式 “2gL sinθ” 中 g 与 L 之间有重叠笔迹（疑有一个 2：V^2=2g·2L sinθ）；图内标注为 “V sinθ 失去”“只剩 V cosθ”
%%%END
%%%BLOCK id=089-02 ch=04 sec=水平面圆周·质心与滑动临界 kind=problem alt=03 cont=none y=0.125-0.255
\begin{problem}
%FIG p089_a 0.10 0.13 0.25 0.205
\notefig[0.25]{figures/p089_a.png}
\begin{solution}
质心 $=\dfrac{2rm+mr}{2m}=\dfrac32r$ 距轴 $\dfrac12r$\qquad $2m\omega^2\dfrac r2>2\mu mg$ 时 向左滑。

若质心距轴为0 则不会滑动\qquad $2m\omega^2\cdot0=0$。
\end{solution}
\end{problem}
%%%END
%%%BLOCK id=089-03 ch=04 sec=圆周运动·圆心参考系 kind=problem alt=07 cont=none y=0.24-0.345
\noindent 圆周运动以圆心为参考系\qquad $V$ 为相对于圆心的速度。
\begin{problem}
%FIG p089_b 0.045 0.27 0.16 0.34
\notefig[0.25]{figures/p089_b.png}
当B\underline{第一次回到$A$}正下方时，求细线拉力大小
% 下划线画在“第一次回到A”下面
\begin{solution}
完全弹性碰撞\qquad $\Delta V_{\text{前}}=\Delta V_{\text{后}}$\qquad $T-mg=\dfrac{m\Delta V^2}{L}=\dfrac{mV_0^2}{L}$
% “ΔV后”的 Δ 手写形似 0
\end{solution}
\end{problem}
%%%END
%%%BLOCK id=090-03 ch=04 sec=平抛·圆轨道出口 kind=problem alt=06 cont=none y=0.33-0.47
\begin{problem}
%FIG p090_c 0.075 0.335 0.285 0.43
\notefig[0.35]{figures/p090_c.png}
求飞得最远的 $R$
\begin{solution}
\[ \begin{cases} S=V_0t\\ 2R=\dfrac12gt^2\\ \dfrac12mV^2-\dfrac12mV_0^2=-mg\,2R\end{cases}\qquad S=\sqrt{4R\left(\dfrac{V_0^2}{g}-4R\right)} \]
% CHECK: 第一式手写形似 S=V_0 t（按物理应为出口速度 V：S=Vt）
均值不等式\quad $R=\dfrac{V_0^2}{8g}$ 时取等
\end{solution}
\end{problem}
%%%END
%%%BLOCK id=090-05 ch=04 sec=斜抛·最大射程 kind=derivation alt=none cont=none y=0.61-0.87
\noindent \textbf{斜抛最大射程}\qquad 法1
%FIG p090_e 0.045 0.632 0.215 0.735
\notefig[0.3]{figures/p090_e.png}
\[ \begin{cases} S=V_0\cos\theta\,t\\ -h=V_0\sin\theta\,t-\dfrac12gt^2\end{cases} \]
\[ -h=V_0\sin\theta\,\frac{S}{V_0\cos\theta}-\frac12g\,\frac{S^2}{(V_0\cos\theta)^2}\qquad\text{由}\ \frac{1}{\cos^2\theta}=1+\tan^2\theta \]
\[ \therefore\ \left(\frac{S}{V_0}\right)^2(1+\tan^2\theta)-\frac{2S}{g}\tan\theta-\frac{2h}{g}=0 \]
\[ \tan^2\theta-\frac{2SV_0^2}{gS^2}\tan\theta+\left(1-\frac{2hV_0^2}{gS^2}\right)=0 \]
\[ \Delta\ge0\ \ \text{故}\ \ S\le\sqrt{\frac{V_0^4}{g^2}+\frac{2V_0^2h}{g}}=\frac{V_0\sqrt{V_0^2+2gh}}{g}=\frac{V_0V_t}{g}\qquad\text{最大射程} \]
\noindent 最大高度\quad $L=\dfrac{(V_0\sin\theta)^2}{2g}$
%%%END
%%%BLOCK id=090-06 ch=04 sec=斜抛·最大射程 kind=derivation alt=none cont=none y=0.86-1.0
\noindent 法2
%FIG p090_f 0.08 0.882 0.185 0.97
\notefig[0.25]{figures/p090_f.png}
\begin{align*}
V_t&=\sqrt{V_0^2+2gh}\\
x&=V_0t\cos\alpha\\
mgt&=m\Delta V\\
\Delta V&=V_0\sin\alpha+V_t\sin\beta\qquad V_0\cos\alpha=V_t\cos\beta
\end{align*}
\[ \therefore\ x=\frac{V_0V_t}{g}\sin(\alpha+\beta) \]
\[ \alpha+\beta=\frac{\pi}{2}\ \text{时}\ x_{\max}=\frac{V_0V_t}{g}\qquad \tan\alpha=\frac{V_0}{V_t} \]
%%%END
%%%BLOCK id=091-04 ch=04 sec=关联速度·绳端速度 kind=derivation alt=01,90 cont=none y=0.33-0.47
%FIG p091_c 0.05 0.33 0.285 0.445
\notefig[0.25]{figures/p091_c.png}

$(h^{2}+x^{2}=l^{2})'$\quad 隐函数求导
\[
2x\frac{\dd x}{\dd t}=2l\frac{\dd l}{\dd t}\qquad \therefore\ v=v_{\text{物}}\frac{x}{l}=v_{\text{物}}\cos\theta
\]
\[
v_{\text{物}}=\frac{v}{\cos\theta}=\frac{\sqrt{x^{2}+h^{2}}}{x}\,v
\qquad
a_{\text{物}}=\frac{\dd v_{\text{物}}}{\dd t}=\left(\frac{\dfrac{1}{2}\cdot\dfrac{2x^{2}v_{\text{物}}}{\sqrt{x^{2}+h^{2}}}-\sqrt{x^{2}+h^{2}}\,v_{\text{物}}}{x^{2}}\right)v
\]
% CHECK: a物 式中分子第一项前的 1/2 手写不太清楚
%%%END
%%%BLOCK id=099-04 ch=04 sec=竖直面匀速圆周运动·洗衣机 kind=problem alt=none cont=none y=0.42-0.54
④
%FIG p099_e 0.055 0.425 0.238 0.528
\notefig[0.25]{figures/p099_e.png}

要衣物\unclear{贴}内壁则，$\omega\ge\sqrt{\dfrac{g}{R}}$

匀速圆周运动

洗衣机
%%%END
%%%BLOCK id=114-01 ch=04 sec=平抛 kind=problem alt=06,90 cont=none y=0.03-0.33
% 本题只有标题“交弧用轨迹方程”、图和“求 Ek min”，没有写出完整题干
\begin{problem}
\textbf{交\unclear{弧}用轨迹方程}

%FIG p114_a 0.01 0.085 0.235 0.19
\notefig[0.3]{figures/p114_a.png}

求 $E_{k\min}$
\begin{solution}
\begin{align*}
y&=-\frac{g}{2v_0^2}x^2+2h=\frac{x^2}{2h}\\
\therefore\ y&=\frac{2v_0^2h}{gh+v_0^2}
\end{align*}
先表达 $x^2=\dfrac{4hv_0^2h}{gh+v_0^2}=\dfrac{4v_0^2h^2}{gh+v_0^2}$

\struck{$y=\unclear{\ldots}$}\quad $y=\dfrac{1}{2h}x^2\ \Rightarrow\ y$．
\begin{align*}
E_k&=mg(2h-y)+\frac{1}{2}mv_0^2\qquad\swarrow\\
&=\frac{2mg^2h^2}{gh+v_0^2}+\frac{1}{2}mv_0^2\\
&=\frac{2mg^2h^2}{gh+v_0^2}+\frac{1}{2}m(gh+v_0^2)-\frac{1}{2}mgh\\
&\ge2\sqrt{2mg^2h^2\cdot\frac{1}{2}m}-\frac{1}{2}mgh=\frac{3}{2}mgh
\end{align*}
\end{solution}
\end{problem}
%%%END
%%%BLOCK id=117-07 ch=04 sec=圆周运动·同轴转动 kind=formula alt=05 cont=none y=0.77-0.96
\begin{align*}
v&=\omega r & v&\propto r\\
a&=\frac{v^2}{r}=\omega^2 r & a&\propto r\\
T&=\frac{2\pi}{\omega}\propto\frac{1}{\omega}\\
\omega&=\frac{2\pi}{T}=2\pi f & f&\propto\omega\\
F_n&=m\omega^2 r\propto m\text{、}r.
\end{align*}
%%%END
%%%BLOCK id=118-01 ch=04 sec=曲线运动·轨迹与合力方向 kind=knowledge alt=09 cont=none y=0.05-0.18
% 页面右上角露出邻页(化学)内容“混合气体 ρ=1.429g/L … 22.4 … 6.72L”，不属于本页，已忽略
\textbf{曲线运动}

速度合力夹轨迹\quad 凹侧方向偏向力

切线$v$\quad 凹侧力\quad 锐加钝减
%%%END
%%%BLOCK id=133-04 ch=04 sec=斜抛·曲率半径 kind=derivation alt=05 cont=none y=0.66-0.84
%FIG p133_d 0.05 0.655 0.455 0.842
\notefig[0.4]{figures/p133_d.png}
求最高点曲率半径
\[ mg=\frac{m\,(V\cos\alpha)^{2}}{\rho} \]
\[ \rho=\frac{V^{2}\cos^{2}\alpha}{g} \] % CHECK: 公式中 V 均未写下标 0（图中初速度为 V_0）
%FIG p133_e 0.85 0.77 1.0 0.84
\notefig[0.25]{figures/p133_e.png}
（页面右侧边缘的小草图，被页边截断：$\rho=\dfrac{v_x^{2}}{\text{\unclear{$a_n$}}}$，旁有 $g$ 字样）
%%%END
%%%BLOCK id=133-05 ch=04 sec=平抛·曲率半径 kind=derivation alt=05 cont=none y=0.84-1.0
%FIG p133_f 0.05 0.841 0.295 0.975
\notefig[0.3]{figures/p133_f.png}
经 $t$ 后 求此时曲率半径
\[ mg\cos\theta=\frac{mV_0^{2}}{\rho\cos\theta} \] % CHECK: 分母 cos 处原稿被涂改（按后一式应为 cos^2θ）
\[ \rho=\frac{mV_0}{\cos^{2}\theta\; mg\cos\theta}=\frac{V_0}{g\cos^{3}\theta}=\frac{V_0}{g\,(g^{2}t^{2}+V_0^{2})^{3/2}} \] % CHECK: 第一个分子 V_0 未见平方；末式 V_0 在分子、(g^2t^2+V_0^2)^{3/2} 在分母，与物理结果（应在分子、V_0 在分母）似乎颠倒，照原稿转录
\[ \tan\theta=\frac{gt}{V_0} \]
%FIG p133_g 0.76 0.838 0.93 0.907
\notefig[0.25]{figures/p133_g.png}
\[ \cos\theta=\frac{V_0}{\sqrt{g^{2}t^{2}+V_0^{2}}} \]
%%%END
%%%BLOCK id=137-02 ch=04 sec=圆周运动·离心运动 kind=knowledge alt=none cont=none y=0.26-0.39
\underline{离转心运}% CHECK: 原稿第2、4字形似“转”“运”，疑为“离心运动”，按原稿转录

$a=\omega^2 r$

$\omega$ 相同 $r$ 越大 $a$ 越大 · $F$\unclear{向}也越大。% CHECK: “向”字原稿形似“方”
%%%END
%%%BLOCK id=164-02 ch=04 sec=圆周运动·竖直面杆模型 kind=note alt=03 cont=none y=0.44-0.56
% 此块上方有一条横贯页面的波浪分隔线(y约0.44)
%FIG p164_d 0.09 0.455 0.25 0.555
\notefig[0.2]{figures/p164_d.png}
对$m$\quad $v$先$\uparrow$后$\downarrow$\quad $a$先右后左.\quad 杆先压缩后拉伸\quad 地先右后左.

$F_{\text{杆}m}=0$时\quad $N=Mg$\quad $f=0$
%%%END
%%%BLOCK id=165-01 ch=04 sec=平抛运动·反弹碰撞 kind=problem alt=06,07 cont=none y=0.07-0.41
\begin{problem}
% 剪贴印刷题干。左缘（第2、4行起首处）有红笔小字/记号，辨认不清；选项B、C左侧有手写记号；“正确”二字被圈出；题干与选项之间有铅笔手绘的折线（V形）划痕
\red{\unclear{…}}甲、乙两同学在课外活动时进行乒乓球比赛，已知乒乓球的质量为 $m$ 且可视为质点，球桌的长度为 $L$，甲同学在球桌左边缘上方高 $H$ 处沿水平方向发球，在桌面上反弹后恰好过网，之后乒乓球落在球桌的右边缘上。乒乓球第一次落在桌面上时的水平位移为 $\frac{3}{8}L$，已知第一次与桌面接触的过程乒乓球与桌面相互作用的时间为 $t$，接触前后水平方向速度不变，重力加速度大小为 $g$，不计摩擦和空气阻力。下列说法正确的是
\begin{enumerate}[label=\Alph*.,leftmargin=2em,itemsep=0pt]
\item 乒乓球第一次与桌面接触过程中损失的能量为 \unclear{$\frac{2}{3}mgH$} % 分数上有一道斜划，分母看不清
\item 乒乓球第一次与桌面接触后反弹上升的最大高度为 $\frac{25}{36}H$
\item 第一次与桌面接触的过程中，乒乓球受到桌面提供的平均作用力为 $\dfrac{11m\sqrt{2gH}}{6t}+mg$
\item 球网的高度为 \struck{$\frac{5}{9}H$} % 印刷选项D被斜线划去
\end{enumerate}
%FIG p165_a 0.36 0.205 0.545 0.30
\notefig[0.3]{figures/p165_a.png}
手写（淡铅笔，图右上方）：BC
\begin{solution}
$V_{y_1}=\sqrt{2gH}$

$V_{y_2}=\sqrt{2gh}$

由动量定理\quad $-\bar{F}t+mgt=-mV_{y_2}-mV_{y_1}$

\[
\bar{F}=\frac{11m\sqrt{2gH}}{6t}+mg
\]

由运动学公式 网高\unclear{$\frac{4}{9}H$} 可建\unclear{系}

两次平抛
\[
\left\{\begin{aligned}
v_0\sqrt{\frac{2H}{g}}&=\frac{3}{8}L\\
2v_0\sqrt{\frac{2h}{g}}&=\frac{5}{8}L
\end{aligned}\right.
\qquad h=\frac{25}{36}L\qquad \Delta E_{\text{损}}=\frac{11}{36}mgH
\]
% CHECK: h=25/36 L 一处原稿末尾写作 L（据上下文应为 H）
\end{solution}
\end{problem}
%%%END
%%%BLOCK id=166-01 ch=04 sec=圆周运动·过弯最短时间 kind=problem alt=none cont=none y=0.00-0.30
\begin{problem}
\textbf{力学也有相切}

车利用技术 用最短时间过 \unclear{路宽} $d=10\unit{m}$ % “d=”写在“路宽”二字上方；“路宽”二字潦草
\red{$r_1=\unclear{70}\unit{m}$\quad $\mu=1.25$}
% CHECK: 红字 r_1 之后的数值字形像“70”（有可能是 7.5 或 10；由后面 r'=12.5m、θ=53° 的 3-4-5 三角形反推 r_1 应为 10 或 7.5），照原样存疑
%FIG p166_a 0.148 0.056 0.362 0.182 soft
\notefig[0.3]{figures/p166_a.png}
\begin{solution}
\red{$r'^2=r_1^2+\left[r'-\left(r_1-\dfrac{d}{2}\right)\right]^2$}

\red{$r'=12.5\unit{m}$}

\red{$\mu mg=\dfrac{mv'^2}{r'}$}

\red{$v'=12.5\unit{m/s}$}

\red{$\theta=53^\circ$}

\red{$t=\dfrac{2\theta}{360^\circ}\cdot\dfrac{2\pi r'}{v'}=1.85\unit{s}$}
\end{solution}
\end{problem}
%%%END
%%%BLOCK id=182-01 ch=04 sec=向心力·传感器实验 kind=method exp=1 alt=none cont=none y=0.00-0.255
% 页顶一行被拍摄裁切，仅见下半截：大致为“拉(住)的力 T=mω…”，其下有分母 r、T² 等残迹，右侧另有零星残笔
\unclear{拉住的力 $T=m\omega\cdots$}

%FIG p182_a 0.065 0.01 0.305 0.118
\notefig[0.3]{figures/p182_a.png}

转 $n$ 圈用时 $t$

$b$烧断前
\[ F_a=mg\qquad \text{测 }m \]
\[ T=\frac{t}{n} \]
\[ r=l_b+\frac{d}{2} \]
\[ F_{\text{向}}=\frac{4\pi^2n^2F_a}{gt^2}\left(l_b+\frac{d}{2}\right) \]
\[ v=\frac{2\pi n}{t}\left(l_b+\frac{d}{2}\right) \]

$b$烧断后\quad 瞬间传感器示数为 $F_0$
\[ F_{\text{向}}=F_0-mg=\frac{F_a}{g}\cdot\frac{v^2}{l_a+\frac{d}{2}}\quad\Rightarrow\ \text{证}\ F_0=F_a\left[\frac{4\pi^2n^2\left(l_b+\frac{d}{2}\right)^2}{gt^2\left(l_a+\frac{d}{2}\right)}+1\right] \]
% “证”字写在 ⇒ 的下方

若$\uparrow$这个等式成立，则向心力公式得到验证
%%%END
%%%BLOCK id=195-03 ch=04 sec=竖直面圆周·平抛临界 kind=problem alt=none cont=none y=0.54-0.88
% 原稿此块（及下方一行）被一条自右上向左下、末端回钩的大弧线划过，疑似划去，仍照录
\notefig[0.4]{figures/p195_b.png}

$R=0.4\unit{m}$

经$e$点$v$越小，下落高度越大，所用时间越长

物体刚好过$e$时速度最小

\begin{align*}
v_e&=\sqrt{gR}=2\unit{m/s}\\
x&=v_e t\\
y&=\tfrac{1}{2}gt^2
\end{align*}

\[
\frac{R-\dfrac{y-R}{\cos\theta}\ \star}{x-(y-R)\tan\theta}=\sin\theta
\]
% 分子右侧有一个小星号标记；分母括号内手写形似“4”，CHECK: 按几何关系应为 $y$

$\therefore\ t=0.3\unit{s}$
%%%END
%%%BLOCK id=207-01 ch=04 sec=水平圆周·转盘连接体临界 kind=method alt=02 cont=none y=0.04-0.26
\textbf{单边型\unclear{水}平圆周}
%FIG p207_a 0.09 0.085 0.285 0.21
\notefig[0.28]{figures/p207_a.png}
$\omega$从0开始$\uparrow$\quad 代入 $f_{\max}=\mu mg$
\begin{itemize}
\item $\omega\in\left[0,\sqrt{\dfrac{\mu g}{2r}}\right)$\quad 都静止\quad $T=0$\quad $f\uparrow$
\item $\omega\in\left(\sqrt{\dfrac{\mu g}{2r}},\sqrt{\dfrac{2\mu g}{3r}}\right)$\quad $f_A$ $T_A$指向圆心，$T_B$\unclear{背}离圆心\quad $f_B$向圆心
\item $\omega\in\left(\sqrt{\dfrac{2\mu g}{3r}},+\infty\right)$\quad $A$ $B$向外滑。（离心运动）
\end{itemize}
\textbf{红笔批注（右侧）：}
\begin{itemize}
\item \red{外侧先打滑。}
\item \red{外侧打滑后，绳为\unclear{外}侧提供拉力与$f$合力充当向心力，对内侧\unclear{……}} % 原稿“绳为?侧”处字被涂改，“对内侧”后有涂抹的字迹
\item \red{$f-T=m\omega^2r_{\text{内}}$}
\item \red{$\omega\uparrow$ $T\uparrow$ $f\uparrow$到$f_{\max}$后，质心打滑}
\end{itemize}
%%%END
%%%BLOCK id=207-02 ch=04 sec=水平圆周·转盘连接体临界 kind=method alt=02 cont=none y=0.26-0.49
\textbf{双边型}
%FIG p207_b 0.09 0.285 0.285 0.445
\notefig[0.28]{figures/p207_b.png}
\begin{itemize}
\item $\omega\in\left[0,\sqrt{\dfrac{\mu g}{2r}}\right]$\quad $f_A$、$f_B$向里\quad $T=0$
\item $\omega\in\left(\sqrt{\dfrac{\mu g}{2r}},\sqrt{\dfrac{\mu g}{r}}\right]$\quad $f_A$、$f_B$向里 有$T$\quad $T$向里
\item $\omega\in\left(\sqrt{\dfrac{\mu g}{r}},\sqrt{\dfrac{2\mu g}{r}}\right]$\quad $f_B$向外，$f_A$向里
\item $\omega\in\left(\sqrt{\dfrac{2\mu g}{r}},+\infty\right)$\quad 质心做离心运动 向左。
\end{itemize}
\textbf{红笔批注（右侧）：}
\begin{itemize}
\item \red{外侧先打滑}
\item \red{\unclear{于是}外侧打滑后}
\item \red{绳为外侧提供拉力}
\item \red{与$f$的合力充当向心力}
\item \red{随$\omega\uparrow$ $T\uparrow$}
\item \red{$f_{\text{内}}\downarrow$到0后反向}
\item \red{增大到$f_{\max}$后}
\item \red{（质心）整体打滑。}
\end{itemize}
%%%END
%%%BLOCK id=207-03 ch=04 sec=水平圆周·转盘连接体临界 kind=method alt=none cont=none y=0.43-0.57
\begin{itemize}
\item $m_Ar_A>m_Br_B$\quad 质心在$A$侧\quad 向$A$侧离心运动
\item $m_Ar_A=m_Br_B$\quad 质心在$O$ 不会打滑，\unclear{张力}$T$一直$\uparrow$直到绳拉断
\item $m_Ar_A<m_Br_B$\quad 质心在$B$侧\quad 向$B$侧离心运动
\end{itemize}
%FIG p207_c 0.49 0.435 0.69 0.56
\notefig[0.3]{figures/p207_c.png}
\[
\mu\,2mg=2m\omega^2\,\frac{r}{2}
\qquad
\omega=\sqrt{\frac{\mu g}{r}}
\]
%%%END
%%%BLOCK id=207-04 ch=04 sec=关联速度·滑轮 kind=method alt=none cont=none y=0.54-0.74
\begin{minipage}[t]{0.40\linewidth}
%FIG p207_d 0.09 0.545 0.205 0.685
\notefig[0.95]{figures/p207_d.png}
$V_B=\frac12V_A$
\end{minipage}\hfill
\begin{minipage}[t]{0.55\linewidth}
%FIG p207_e 0.28 0.545 0.49 0.685
\notefig[0.95]{figures/p207_e.png}
$V_B=\frac12V_A\cos\alpha$
\end{minipage}

\red{竖直速度为2倍关系}\quad \fbox{$V_A\cos\alpha=2V_B$}
%%%END
%%%BLOCK id=208-02 ch=04 sec=平抛·斜面上的平抛与等时间间隔投弹 kind=diagram alt=none cont=none y=0.58-0.72
\begin{minipage}[t]{0.47\linewidth}
%FIG p208_b 0.04 0.585 0.35 0.72
\notefig[0.95]{figures/p208_b.png}
\end{minipage}\hfill
\begin{minipage}[t]{0.50\linewidth}
等时间间隔投弹\quad $t_{ab}=2t_{bc}$
%FIG p208_c 0.40 0.607 0.73 0.725
\notefig[0.95]{figures/p208_c.png}
\end{minipage}

图中文字：左图“更早到”；右图“第3个在$bc$之间”，底线下方“第一个 第2个”。
%%%END
%%%BLOCK id=208-03 ch=04 sec=平抛·斜面上的平抛与等时间间隔投弹 kind=method alt=none cont=none y=0.73-0.95
\begin{minipage}[t]{0.47\linewidth}
%FIG p208_d 0.04 0.745 0.32 0.86
\notefig[0.95]{figures/p208_d.png}
由 $x_1=x_2$ $\therefore\Delta t_1=\Delta t_2$
\end{minipage}\hfill
\begin{minipage}[t]{0.50\linewidth}
%FIG p208_e 0.45 0.745 0.75 0.86
\notefig[0.95]{figures/p208_e.png}
$x_2=2x_1$
\end{minipage}

右图中文字：飞机；以$V_0$射出；以$2V_0$射出。

\unclear{定}同一地面\quad 水平\quad $x\propto v$
%%%END
%%%BLOCK id=209-01 ch=04 sec=曲线运动·曲率半径与法向力 kind=method alt=11 cont=none y=0.10-0.56
洛伦兹力下 $R\uparrow$ $V\uparrow$

变速曲线下 $R\uparrow$ $V\downarrow$\qquad $F_n=\dfrac{mV^2}{\rho}$
%FIG p209_a 0.04 0.29 0.26 0.50
\notefig[0.3]{figures/p209_a.png}
$F_n=mg\cos\theta$
\[
\begin{cases}
mg\cos\theta=\dfrac{mV^2}{\rho}\\[4pt]
mgy=\dfrac12mV^2\\[4pt]
\tan\theta=2\tan\alpha=\dfrac{2y}{x}
\end{cases}
\qquad
y=\dfrac{g^2x^2}{2V_0^2}
\]
% CHECK: 右侧 y 的表达式分子写作 g²x²，与平抛轨迹 y=gx²/(2v0²) 不一致，请核对
%%%END
%%%BLOCK id=221-08 ch=04 sec=平抛运动·活用公式 kind=method alt=none cont=none y=0.62-0.77
\textbf{7、}活用公式
%FIG p221_e 0.04 0.64 0.25 0.765
\notefig[0.25]{figures/p221_e.png}
\begin{align*}
t_1&=\sqrt{\frac{2(h_1+h)}{g}}\\
t_2&=\sqrt{\frac{2(h_1+2h)}{g}}
\end{align*}
\[
\frac{t_1}{t_2}=\frac{\sqrt{h_1+h}}{\sqrt{h_1+2h}}<\frac{\sqrt2}{2}\ \unclear{…}\quad\text{故}\ t_2<\sqrt2\,t_1
\] % CHECK: 按算式 t1/t2 应不小于 √2/2，原稿不等号为“<”；“√2/2”后有一处涂改字迹
%%%END
%%%BLOCK id=222-06 ch=04 sec=圆周运动·向心力（相对速度） kind=method alt=07,06 cont=none y=0.64-0.83
向心力。相对于圆心的速度\quad $F_n=\dfrac{mv^2}{r}$

%FIG p222_d 0.0 0.705 0.15 0.785
\notefig[0.25]{figures/p222_d.png}
\begin{align*}
mgR&=\frac12MV'^2+\frac12mV^2\\
mV&=MV'\\
N-mg&=\frac{m(V+V')^2}{r}
\end{align*}
%%%END
%%%BLOCK id=225-03 ch=04 sec=平抛运动·频闪照片 kind=problem exp=1 alt=01 cont=none y=0.15-0.38
\begin{problem}
%FIG p225_b 0.462 0.165 0.74 0.274
\notefig[0.3]{figures/p225_b.png}
$T=0.1\unit{s}$
\begin{solution}
$\Delta y=y_2-y_1=gT^2$\quad $\therefore\ V_0=2\unit{m/s}$

$V_{Ay}=1.5\unit{m/s}$\quad $t=0.15\unit{s}$

$V_A=2.5\unit{m/s}$
\end{solution}
\end{problem}
%%%END
%%%BLOCK id=225-06 ch=04 sec=平抛运动·动能与动量变化 kind=problem alt=06,07 cont=none y=0.37-0.50
\begin{problem}
%FIG p225_c 0.512 0.378 0.600 0.43
\notefig[0.25]{figures/p225_c.png}
\begin{solution}
第3s内：第5s内\quad $E_k$变化\quad $5:9$ % 原稿：“E_k变化”写在第二行，并以括号与 5:9 相连

动量变化之比\quad $1:1$
\end{solution}
\end{problem}
%%%END
%%%BLOCK id=226-09 ch=04 sec=平抛·相对位移 kind=note alt=03,06 cont=none y=0.69-0.78
%FIG p226_g 0.095 0.69 0.165 0.75
\notefig[0.25]{figures/p226_g.png}

$m$ 掉后，同向平抛$x_{\text{相}}=x_M-x_m$

$m$ 掉下后 车为保持匀速 $f=\mu Mg\downarrow$，$F_{\text{牵}}=f$ 也变小
%%%END
%%%BLOCK id=227-01 ch=04 sec=圆周运动·分离临界与关联速度 kind=problem alt=03,06 cont=none y=0.01-0.25
\textbf{1.} 加速关联
%FIG p227_a 0.19 0.005 0.505 0.08
\notefig[0.31]{figures/p227_a.png}

\unclear{将}要分离

$a_{Ax}=a_{Bx}=0$ \qquad $V_A=V_B\sin30^\circ$ % CHECK: 与后面 V_B=√(gL/8) 对照，疑应为 V_B=V_A sin30°

$\therefore a_{Ay}=g$ \qquad $mg\sin30^\circ=\dfrac{mV_A^{2}}{L}$ 重力分力提供向心力

$V_A=\sqrt{\dfrac{3gL}{2}}$ \qquad $m_A:m_B=1:4$ % CHECK: 下标可能是"A末"（末速度）；与上式按 mg sin30°=mV_A^2/L 得 V_A=√(gL/2) 不同

$m_Agl\sin30^\circ=\dfrac12 mV_{A\text{末}}^{2}-\dfrac12 mV_A^{2}$ % CHECK: 末 字辨认不确定

$m_Agl(1-\sin30^\circ)=\dfrac12 m_AV_A^{2}+\dfrac12 m_BV_B^{2}$ \qquad $V_B=\sqrt{\dfrac{gL}{8}}$
%%%END
%%%BLOCK id=227-04 ch=04 sec=竖直面圆周·不脱轨条件 kind=note alt=06 cont=none y=0.30-0.46
\textbf{2.} 不出轨道
%FIG p227_b 0.04 0.35 0.20 0.45
\notefig[0.25]{figures/p227_b.png}

$V\ge V_2$ 或 $V\le V_1$ \qquad $\ast$ 反回或冲出后碰墙反弹回来。
%%%END
%%%BLOCK id=227-05 ch=04 sec=皮带轮·抛出临界 kind=note alt=06 cont=none y=0.46-0.54
%FIG p227_c 0.05 0.47 0.205 0.52
\notefig[0.3]{figures/p227_c.png}

能在D抛出 \quad $mg\le\dfrac{mV_D^{2}}{R}$

求$h$、$x$关系 分段讨论

过了某个点/未过某个点
%%%END
%%%BLOCK id=227-06 ch=04 sec=圆周运动·脱离轨道 kind=derivation alt=06 cont=none y=0.52-0.67
%FIG p227_d 0.04 0.52 0.225 0.64
\notefig[0.3]{figures/p227_d.png}

重力分力提供向心力

在P点 $mg\cos\theta-N=\dfrac{mV_P^{2}}{R}$ \qquad 从B到P $mgR(1-\cos\theta)=\dfrac12 mV_P^{2}$

在P脱离轨道则 $N=0$

$\cos\theta=\dfrac{h}{R}$
%%%END
%%%BLOCK id=227-07 ch=04 sec=水平圆周·圆锥内物块 kind=problem alt=02 cont=none y=0.69-0.87
\textbf{3.} 圆锥
%FIG p227_e 0.04 0.729 0.19 0.84
\notefig[0.25]{figures/p227_e.png}

①A静止时求$f$、$N$ \quad 转动时求$\omega$（$\mu=0$）

\[ f=mg\sin\theta=\frac{mgH}{\sqrt{H^{2}+R^{2}}}\qquad
\begin{cases} mg\tan\theta=m\omega^{2}\dfrac{R}{2}\\[4pt] \tan\theta=\dfrac{H}{R}\end{cases} \]
\[ N=mg\cos\theta=\frac{mgR}{\sqrt{H^{2}+R^{2}}}\qquad \therefore\omega=\frac{\sqrt{2gH}}{R} \]
%%%END
%%%BLOCK id=229-04 ch=04 sec=斜抛·斜面上的抛体 kind=derivation alt=03 cont=none y=0.82-1.00
\textbf{3.}
%FIG p229_c 0.00 0.82 0.195 0.965
\notefig[0.25]{figures/p229_c.png}

\[ t=\frac{2V_{0y}}{g_y}=\frac{2V_0\cos(\alpha-\theta)}{g\cos\theta} \]
也可按斜抛水平竖直分解

当 $\cos(\alpha-\theta)=1$ 时（$\alpha=\theta$）$t$最大

\[ t=\frac{2V_0}{g\cos\theta}\qquad x=\frac12 g_xt^{2}=\frac{2V_0^{2}\sin\theta}{g\cos^{2}\theta} \]
$g_x=g\sin\theta$
%%%END
%%%BLOCK id=235-01 ch=04 sec=竖直面圆周·最低点受力 kind=note alt=06 cont=none y=0.02-0.12
%FIG p235_a 0.025 0.035 0.16 0.108 soft
\notefig[0.2]{figures/p235_a.png}
$1\to2$ 在$2$处 $N=3mg$\par
%FIG p235_b 0.43 0.015 0.54 0.105 soft
\notefig[0.2]{figures/p235_b.png}
$T=3mg$
%%%END
%%%BLOCK id=240-08 ch=04 sec=圆周运动·圆锥摆 kind=knowledge alt=none cont=none y=0.80-0.95
\textbf{圆锥摆}

%FIG p240_i 0.15 0.822 0.275 0.91
\notefig[0.2]{figures/p240_i.png}

多体摆　顶点不重合，全在同一高度

\[
\left\{\begin{array}{l}
T\cos\theta=m\omega^2R\\
T\sin\theta=mg
\end{array}\right.
\qquad
\omega=\sqrt{\frac{g}{h}}
\]
%%%END
%%%BLOCK id=240-09 ch=04 sec=圆周运动·临界问题 kind=method alt=none cont=none y=0.78-0.95
常规方法

两绳上都有拉力。

%FIG p240_j 0.685 0.843 0.828 0.945
\notefig[0.25]{figures/p240_j.png}

\[
\omega_{\max}=\sqrt{\dfrac{g}{\dfrac{\sqrt3}{3}\,r}}
\qquad
\omega_{\min}=\sqrt{\dfrac{g}{\sqrt3\,r}}
\]
%%%END
%%%BLOCK id=240-10 ch=04 sec=圆周运动·临界问题 kind=note alt=02 cont=none y=0.93-0.96
结合$f$的两向性考查
%%%END
%%%BLOCK id=242-03 ch=04 sec=平抛·斜面上的平抛 kind=problem alt=none cont=none y=0.36-0.56
\begin{problem}[3、平抛]
%FIG p242_e 0.075 0.386 0.375 0.507
\notefig[0.3]{figures/p242_e.png}
速度反向延长线\\交于位移中点
\begin{solution}
\begin{align*}
&t_A<t_B<t_C\\
&x_A=x_B=x_C\\
&\therefore V_{A0}>V_{B0}>V_{C0}
\end{align*}
\[ h=\frac{x}{2}\,\frac{1}{\tan\theta}\qquad t=\sqrt{\frac{2h}{g}}\qquad V_y=\sqrt{2gh}\qquad V_x=\tan\theta\sqrt{2gh} \]
\[ \therefore V=\sqrt{V_x^2+V_y^2}=\sqrt{2gh\,(1+\tan^2\theta)}=\sqrt{xg\Bigl(\frac{1}{\tan\theta}+\frac{1}{\tan\theta}\Bigr)} \] % CHECK: 按推导括号内第二项应为 tanθ，原稿写作 1/tanθ
代入 $60^\circ$ $45^\circ$ $30^\circ$ 得 $V_A=V_C>V_B$
\end{solution}
\end{problem}
%%%END
%%%BLOCK id=243-01 ch=04 sec=圆周运动·圆锥摆(双球共轴) kind=problem alt=02 cont=none y=0.03-0.22
\begin{problem}
%FIG p243_a 0.115 0.035 0.275 0.125
\notefig[0.25]{figures/p243_a.png}
\begin{solution}
\begin{align*}
F\cos\alpha&=m_Ag\\
F\cos\beta&=m_Bg\\
m_Ag\tan\alpha&=m_A\omega^2l_1\sin\alpha\\
l_1\cos\unclear{\beta}&=l_2\cos\alpha\\ % CHECK: 第一个 cos 后的角度被涂改（β 与 α 叠写）；第二个像 α；按推导应为 l1cosα=l2cosβ
\frac{m_A}{m_B}&=\frac{\cos\alpha}{\cos\beta}
\end{align*}
\end{solution}
\end{problem}
%%%END
%%%BLOCK id=243-03 ch=04 sec=圆周运动·洗衣机转筒 kind=problem alt=06 cont=none y=0.29-0.43
\begin{problem}
洗衣机转筒
%FIG p243_c 0.09 0.32 0.265 0.405
\notefig[0.25]{figures/p243_c.png}
%FIG p243_d 0.78 0.285 0.925 0.39
\notefig[0.25]{figures/p243_d.png}
\begin{solution}
\[ \begin{cases}
\mu mg=m\omega^2r\\
-\mu mgx=0-\frac12m(\omega x)^2 \\ % CHECK: 括号内字母形似 x（物理上应为 r）
R=\sqrt{x^2+r^2}\qquad \unclear{F}_N>0 \\ % CHECK: 下标 N 前的字母形似 I（也可能是 F）
Q=\mu mg\,\unclear{…}\,x % g 与 x 之间有一个被涂黑的符号
\end{cases} \]
\end{solution}
\end{problem}
%%%END
%%%BLOCK id=243-04 ch=04 sec=圆周运动·合外力与向心力 kind=note alt=none cont=none y=0.44-0.47
单摆秋千合外力不恒\unclear{指向}圆心
%%%END
%%%BLOCK id=243-05 ch=04 sec=平抛·斜面上的平抛(速度偏角) kind=problem alt=none cont=none y=0.48-0.60
\begin{problem}
%FIG p243_e 0.065 0.51 0.235 0.60
\notefig[0.25]{figures/p243_e.png}
\begin{solution}
\[ \tan\theta=\tan(\alpha-\varphi)=\frac{\tan\alpha-\tan\varphi}{\tan\alpha\tan\varphi+1}\qquad\tan\alpha=2\tan\varphi \]
\[ =\frac{1}{2\tan\varphi+\dfrac{1}{\tan\varphi}} \]
%FIG p243_f 0.60 0.52 0.80 0.60
\notefig[0.25]{figures/p243_f.png}
图旁标注：不单调
\end{solution}
\end{problem}
%%%END
%%%BLOCK id=243-06 ch=04 sec=斜抛·落点速度方向 kind=note alt=none cont=none y=0.60-0.64
斜抛 则落点速度方向不同\quad $t\neq\dfrac{2V_0\tan\varphi}{g}$\quad $\Delta y\neq g\Delta t^2$
%%%END
%%%BLOCK id=247-03 ch=04 sec=运动合成与分解·关联速度 kind=derivation alt=none cont=none y=0.59-0.78
%FIG p247_d 0.07 0.595 0.325 0.718
\notefig[0.3]{figures/p247_d.png}
\[
V_{\perp}=V\cos\theta=\omega r\qquad r=\frac{h}{\cos\theta}
\]
% 原稿此式被划去
\struck{$V_{\text{速}}=V_{\text{影}}\cos\theta$}
\begin{keybox}
\[
\sin\text{光地}=\frac{V_{\text{\unclear{顶}}}}{V_{\text{影}}}
\]
\end{keybox}
%%%END
%%%BLOCK id=250-01 ch=04 sec=圆周运动·竖直面 kind=problem alt=06 cont=none y=0.00-0.17
%FIG p250_a 0.05 0.0 0.15 0.052 soft
\notefig[0.25]{figures/p250_a.png}
\[
mg(\sin37)^\circ=\frac{mV'^2}{R}
\]
\[
V'^2=\frac35gR
\]
\[
-mg\,\frac85R=\frac12m\left(\frac35gR-V_0^2\right)
\]
\[
V_0^2=\frac{19}{5}gR\qquad V_0=\sqrt{\frac{19}{5}gR}
\]
%%%END
%%%BLOCK id=250-02 ch=04 sec=圆周运动·竖直面 kind=knowledge alt=none cont=none y=0.01-0.12
不脱轨\quad $V_0\le\sqrt{2gR}$\quad $V\ge\sqrt{5gR}$ \\
脱轨\ \unclear{$\Rightarrow$}\ $V_0\in\left(\sqrt{2gR},\sqrt{5gR}\right)$
%%%END
%%%BLOCK id=250-03 ch=04 sec=圆周运动·竖直面 kind=problem alt=06 cont=none y=0.00-0.16
%FIG p250_b 0.57 0.0 0.725 0.12
\notefig[0.3]{figures/p250_b.png}
$v=\sqrt{2gl}$

使其绕 $P$ \unclear{的}完整\unclear{圆周}运动
\[
\sqrt{2gl}\ge\sqrt{5g(l-d)}
\]
%%%END
%%%BLOCK id=250-05 ch=04 sec=圆周运动·水平面 kind=knowledge alt=none cont=none y=0.28-0.47
\textbf{圆锥摆}
%FIG p250_c 0.065 0.325 0.225 0.468
\notefig[0.25]{figures/p250_c.png}
\[
F_n=mg\tan\theta=\frac{mv^2}{r}=m\omega^2r
\]
\[
r=L\sin\theta\qquad v=\sqrt{gL\sin\theta\tan\theta}\qquad \omega=\sqrt{\frac{g}{L\cos\theta}}=\sqrt{\frac{g}{h}}
\]
若 $\theta\uparrow$ 则 $\omega\uparrow$ $v\uparrow$ $F_n\uparrow$ $T\uparrow$
%%%END
%%%BLOCK id=250-06 ch=04 sec=圆周运动·水平面 kind=knowledge alt=none cont=none y=0.50-0.65
\textbf{圆锥筒}
%FIG p250_d 0.08 0.52 0.30 0.655
\notefig[0.25]{figures/p250_d.png}
\[
\frac{mg}{\tan\theta}=\frac{mv^2}{r}=m\omega^2r
\]
\[
V=\sqrt{\frac{gr}{\tan\theta}}\qquad \omega=\sqrt{\frac{g}{r\tan\theta}}
\]
小球位置越高 $r\uparrow$ $\omega\downarrow$ $v\uparrow$
%%%END
%%%BLOCK id=250-07 ch=04 sec=平抛运动·斜面 kind=knowledge alt=none cont=none y=0.66-0.88
\textbf{平抛斜面}
%FIG p250_e 0.08 0.692 0.285 0.79
\notefig[0.25]{figures/p250_e.png}
位偏角 $\theta$\quad $t=\dfrac{2V_0\tan\theta}{g}$

同角落在斜面 只与 $\theta$ 有关

落回斜面上水平位移与初速度平方成正比

最高点\quad $\dfrac{V_y}{V_x}=\tan\theta$\quad 该刻为全过程中间时刻\quad $t=\dfrac{V_0\tan\theta}{g}$

该点到斜面的最远距离\ $d=\dfrac{V_0^2\sin\theta}{2g\cos\theta}$ % CHECK: 按推导分子似应为 V_0^2 sin^2θ，原稿写 sinθ
%%%END
%%%BLOCK id=250-08 ch=04 sec=斜抛运动 kind=formula alt=none cont=none y=0.88-1.00
\textbf{斜抛运动}
%FIG p250_f 0.05 0.897 0.262 0.97
\notefig[0.25]{figures/p250_f.png}
速度：$V_x=V_{0x}=V_0\cos\theta$\qquad $V_y=V_0\sin\theta=gt$。 % CHECK: 第二个等号原稿为“=”，按物理似应为 V_y=V_0\sin\theta-gt

斜抛运动具有轨迹、速度、运动时间对称，加速度

位移\ $x=V_0\cos\theta\,t$\qquad $y=V_0\sin\theta\,t-\dfrac12gt^2$

射高\ $H=\dfrac{V_{0y}^2}{2g}=\dfrac{(V_0\sin\theta)^2}{2g}$\qquad $t=\dfrac{2V_0\sin\theta}{g}$

水平射程\ $x=V_{0x}t=V_0\cos\theta\,\dfrac{\unclear{2}V_0\sin\theta}{g}$ % 原稿分子前为一小点，按推导应为 2
\[
=\frac{V_0^2\sin2\theta}{g}
\]
%%%END
%%%BLOCK id=258-01 ch=04 sec=圆周运动·基本量 kind=formula alt=none cont=none y=0.07-0.22
\begin{gather*}
v=\frac{\Delta\widehat{l}}{\Delta t}\qquad \unit{m/s}\\
\omega=\frac{\Delta\theta}{\Delta t}\ \ \text{矢量}\qquad \unit{rad/s}\\
n=f=\frac{\omega}{2\pi}\qquad \unit{r/s}
\end{gather*}
匀圆\quad $a_{\text{合}}=a_n$\quad $F_{\text{合}}=F_n$
\[ F_n=\frac{mv^2}{r}=m\omega^2 r=mv\omega=ma=m\frac{4\pi^2}{T^2}r \]
%%%END
%%%BLOCK id=258-02 ch=04 sec=皮带齿轮·传动装置 kind=knowledge alt=none cont=none y=0.17-0.40
\textbf{传动装置}：线速度相等
\begin{itemize}
\item 皮带（静摩擦力）
\item 齿轮传动\quad $T_{\text{分}}=60T_{\text{秒}}$\quad \unclear{$T_{\text{分}}=60\,r_{\text{秒}}$针$+$}
\item $f_{B\text{带}}$\ 链条
\end{itemize}
自行车\ 主\struck{\unclear{…}}动轮\quad 受地\unclear{$f$}向前
%FIG p258_a 0.565 0.225 0.762 0.295
\notefig[0.3]{figures/p258_a.png}
\[ r_A=\frac{1}{2}r_B,\qquad V_A=V_B,\qquad \frac{\omega_A}{\omega_B}=2 \]
受主动轮\unclear{转}方向\par
与主动轮转向一样
%FIG p258_b 0.585 0.300 0.755 0.350
\notefig[0.3]{figures/p258_b.png}
%%%END
%%%BLOCK id=258-03 ch=04 sec=皮带齿轮·传动装置 kind=knowledge alt=none cont=none y=0.37-0.52
\textbf{角速度相等}\ —\ 同轴传动
%FIG p258_c 0.12 0.405 0.285 0.472
\notefig[0.25]{figures/p258_c.png}
\[ r_B=3r_A,\qquad V_B=3V_A \]
%%%END
%%%BLOCK id=258-04 ch=04 sec=皮带齿轮·传动装置 kind=knowledge alt=none cont=none y=0.38-0.54
%FIG p258_d 0.355 0.390 0.762 0.535
\notefig[0.7]{figures/p258_d.png}
\[ V_1=V_2,\qquad \omega r\ \text{相等},\qquad n_1D_1=n_2D_2 \]
可动轴；主；越转越快
\[ r_1\omega_1=r_2\omega_2 \]
主；从；可变半径
%%%END
%%%BLOCK id=259-04 ch=04 sec=圆周运动·匀速圆周 kind=note alt=06 cont=none y=0.53-0.56
匀圆不用动能定理
%%%END
%%%BLOCK id=259-05 ch=04 sec=平抛·斜抛 kind=knowledge alt=none cont=none y=0.60-0.77
\textbf{斜抛结论}
%FIG p259_d 0.075 0.617 0.230 0.770
\notefig[0.25]{figures/p259_d.png}
斜抛在 $\alpha$ 角平分线处\ 水平距离最大
%%%END
%%%BLOCK id=260-01 ch=04 sec=皮带齿轮·传动装置 kind=summary alt=none cont=none y=0.00-0.30
\textbf{皮带传动}
%FIG p260_a 0.110 0.026 0.255 0.095
\notefig[0.25]{figures/p260_a.png}
\[ V_A=V_C\neq V_B \]
%FIG p260_b 0.430 0.005 0.690 0.115
\notefig[0.3]{figures/p260_b.png}
\[ V_A=V_B\neq V_P=V_Q \]
\textbf{摩擦传动、齿轮传动}
%FIG p260_c 0.090 0.126 0.225 0.166
\notefig[0.25]{figures/p260_c.png}
\[ V_A=V_B \]
\textbf{同轴传动}
%FIG p260_d 0.100 0.191 0.205 0.265
\notefig[0.25]{figures/p260_d.png}
%FIG p260_e 0.250 0.168 0.360 0.275
\notefig[0.25]{figures/p260_e.png}
\[ \omega_A=\omega_B,\qquad V\propto r \]
%%%END
%%%BLOCK id=260-02 ch=04 sec=水平面圆周·向心力来源 kind=diagram alt=none cont=none y=0.33-0.45
水平飞机转弯
%FIG p260_f 0.045 0.335 0.190 0.450
\notefig[0.25]{figures/p260_f.png}
火车转弯
%FIG p260_g 0.200 0.335 0.480 0.445
\notefig[0.4]{figures/p260_g.png}
\[ F_n=mg\tan\theta \]
%%%END
%%%BLOCK id=260-03 ch=04 sec=水平面圆周·向心力来源 kind=diagram alt=none cont=none y=0.455-0.575
%FIG p260_h 0.040 0.455 0.190 0.575
\notefig[0.25]{figures/p260_h.png}
\[ F_n=mg\tan\theta,\qquad r=l\sin\theta \]
%FIG p260_i 0.325 0.485 0.520 0.570
\notefig[0.3]{figures/p260_i.png}
\[ F_n=mg\tan\theta \]
%%%END
%%%BLOCK id=260-04 ch=04 sec=水平面圆周·向心力来源 kind=diagram alt=none cont=none y=0.585-0.69
%FIG p260_j 0.070 0.595 0.170 0.675
\notefig[0.25]{figures/p260_j.png}
\[ F_n=f \]
%FIG p260_k 0.275 0.583 0.520 0.695
\notefig[0.3]{figures/p260_k.png}
（光滑杆\ 转台）、\par
\[ F_n=m_Bg \]
%%%END
%%%BLOCK id=260-05 ch=04 sec=水平面圆周·圆锥摆 kind=derivation alt=none cont=none y=0.71-0.95
\textbf{圆锥摆}
%FIG p260_l 0.000 0.735 0.170 0.855
\notefig[0.25]{figures/p260_l.png}
\[ mg\tan\theta=m\omega^2L\sin\theta=m\frac{4\pi^2}{T^2}L\sin\theta \]
\begin{align*}
\omega&=\sqrt{\frac{g}{L\cos\theta}} & T&=2\pi\sqrt{\frac{L\cos\theta}{g}}\\
&=\sqrt{\frac{g}{h}}
\end{align*}
\[ T=m\omega^2L \]
%%%END
%%%BLOCK id=260-06 ch=04 sec=水平面圆周·圆锥摆 kind=method alt=none cont=none y=0.65-1.00
%FIG p260_m 0.486 0.650 0.880 0.847
\notefig[0.45]{figures/p260_m.png}
同高不同顶\par
等效绳长。
\[ \omega_1=\omega_2 \]
\[ V\propto r,\qquad a\propto r \]
竖直方向合力为0
\[
\left\{\begin{aligned}
r'&=r+L\sin\alpha\quad\text{等效绳长}\\
F\sin\alpha&=m\omega^2(r+L\sin\alpha)\\
F&=m\left(\frac{r}{\sin\alpha}+L\right)\omega^2=\frac{mg}{\cos\alpha}\\
\omega&=\sqrt{\frac{g}{r'\cos\alpha}}\qquad T=2\pi\sqrt{\frac{r'\cos\alpha}{g}}
\end{aligned}\right.
\] % CHECK: 末行 r' 的含义与首行不一致（由上式推得应为 ω^2=g tanα/r'），原样转录
%%%END
%%%BLOCK id=261-01 ch=04 sec=水平面圆周·飞行器综合题 kind=problem alt=02,06 cont=none y=0.00-0.93
\begin{problem}
% 原稿：粘贴的印刷题，首行从「到影响，」起（上文在粘贴时被裁掉）。印刷文字上有手画下划线/圈注：「飞行器的动力F始终与飞行方向相同」「空气升力F1与飞行方向垂直」「大小与速度的平方成正比，即F1=C1v^2」「大小与速度的平方成正比，即F2=C2v^2」「可由运动员调节」下划线；「总质量为90kg」被圈出；「（重力加速度取g=10m/s^2。）」上方有弧线。左页边有零星手写标记。
到影响，同时通过控制动力的大小而改变飞行器的飞行状态。已知，飞行器的动力 $F$ 始终与飞行方向相同，空气升力 $F_1$ 与飞行方向垂直，大小与速度的平方成正比，即 $F_1=C_1v^2$；空气阻力 $F_2$ 与飞行方向相反，大小与速度的平方成正比，即 $F_2=C_2v^2$，其中 $C_1$、$C_2$ 相互影响，可由运动员调节，满足如图1所示的关系。飞行员和装备的总质量为 $90\unit{kg}$。（重力加速度取 $g=10\unit{m/s^2}$。）
\begin{enumerate}[label=（\arabic*）]
\item 若飞行员使飞行器以 $v_1=10\sqrt{3}\unit{m/s}$ 速度在空中沿水平方向匀速飞行，如图2所示，则飞行器受到的动力 $F$ 大小为多少？
\item 若飞行员关闭飞行器的动力，使飞行器匀速滑行，且滑行速度 $v_2$ 与地平线的夹角 $\theta=30^\circ$，如图3所示，则速度 $v_2$ 的大小为多少？（结果可用根式表示）
\item 若飞行员使飞行器在空中的某一水平面内做匀速圆周运动，如图4所示，在此过程中 $C_2$ 只能在 $1.75\sim2.5\unit{N\cdot s^2/m^2}$ 之间调节，且 $C_1$、$C_2$ 的大小与飞行器的倾斜程度无关。则飞行器绕行一周动力 $F$ 做功的最小值为多少？（结果可保留 $\pi$）
\end{enumerate}
%FIG p261_a 0.015 0.380 0.400 0.595
\notefig[0.4]{figures/p261_a.png}
% 图1上手画：从原点向右上的细斜线，与曲线交于 C_1=4 处，并由交点向下画竖线至横轴，横轴下方手写「2.3」
\red{（$-7$）}\quad\red{$-2$} % 红笔批注：前者被红圈圈住，在图2~4左上方空白处；后者在图1右侧
%FIG p261_b 0.405 0.355 1.000 0.580
\notefig[0.8]{figures/p261_b.png}
% 图3上手画受力箭头并标 F_2、F_1、mg；图4上手画上下箭头
% 图4上方/右侧手写草算（右侧被页边截断）：
\par $C_2$\,\struck{$v^2$}$\,=\dfrac{900}{r}$\quad\struck{$v^2$}\quad $C_2=\dfrac{\unclear{90}}{r}$
\par $r=\dfrac{\unclear{90}}{C_2}$\quad $=\dfrac{\unclear{0}}{\unclear{\ldots}}$\quad $=\unclear{36}$
\par 图4下方：$\dfrac{36}{\unclear{\ldots}\tan\unclear{\alpha}}$
\begin{solution}
(1)\quad $\overset{\text{平衡}}{mg=C_1v_1^2}$\quad $C_1=3\unit{N\cdot s^2/m^2}$
\par $\therefore C_2=2.5\unit{N\cdot s^2/m^2}$
\par 动力 $F=F_2=C_2v_1^2=750\unit{N}$
\par (2)\quad $mg\cos\theta\overset{\text{平衡}}{=}C_1v_2^2$
\par $mg\sin\theta=C_2v_2^2$
\par $C_1=V_2\cot\theta$ % CHECK: 两式相除应得 C_1=C_2\cot\theta，原稿右端写作 V_2（或为 C_2 写得像 V）
\par $C_2=2.3\unit{N\cdot s^2/m^2}$\quad $C_1=4\unit{N\cdot s^2/m^2}$
\par $V_2=\sqrt{\dfrac{225\sqrt{3}}{2}}\unit{m/s}$
\par (3)\quad 非平衡
\par 设 $F_1$ 与竖直方向夹角 $\alpha$
\par $mg=C_1v^2\cos\alpha$
\par 水平合力 提供向心力\quad $C_1v^2\sin\alpha=\dfrac{mv^2}{R}$
\par 动力\quad $F=F_2=C_2v^2$
\par $W=F\cdot 2\pi R=\dfrac{2\pi C_2m^2g}{C_1^2\sin\alpha\cos\alpha}$
\par $\alpha=45^\circ$\unclear{，}\ $C_1=6\unit{N\cdot s^2/m^2}$，$C_2=1.75\unit{N\cdot s^2/m^2}$\ 时 $W$ 最小， % 原稿折行：「C_2=1.75…」写在右侧，「时 W最小，」写在左侧下一行；「45°」后有一个难辨的小字
\par $W_{\min}=15750\pi\unit{J}$
\end{solution}
\end{problem}
%%%END
%%%BLOCK id=262-01 ch=04 sec=圆周运动·解题步骤 kind=summary alt=none cont=none y=0.04-0.20
\textbf{圆周运动}
\par 一、确定轨道平面、圆心位置、轨迹圆半径
\par 二、受力分析、找向心力、正交分解、指向圆心减背离圆心的力
\par 三、列方程求解
\begin{itemize}
\item 点、牛二
\item 过程、动能定理
\end{itemize}
%%%END
%%%BLOCK id=262-02 ch=04 sec=竖直面圆周·绳模型 kind=knowledge alt=none cont=none y=0.21-0.47
\textbf{绳模型}
\par 绳只能拉不能推，圆心以上\ 速度一定不为0
%FIG p262_a 0.185 0.272 0.253 0.325
\notefig[0.25]{figures/p262_a.png}
\par 最高点\quad $F+G=\dfrac{mv^2}{L}$\qquad $F\ge0$\quad $\therefore V\ge\sqrt{gL}$
\par 最低点\quad $\left\{\begin{aligned}F-G&=\dfrac{mv^2}{L}\\ mg\cdot2L&=\dfrac{1}{2}m(v_0^2-v^2)\end{aligned}\right.$\qquad $\begin{aligned}V^2&=V_0^2-4gl\ge gl\\ \unclear{\therefore}\,V_0&\ge\sqrt{5gL}\end{aligned}$
\par 拉力差\quad $\Delta F=6mg$
\begin{align*}
mg+F_1&=\frac{mv^2}{L}\\
F_2-mg&=\frac{mv'^2}{L}\\
2mgl&=\frac{1}{2}m(v'^2-v^2)
\end{align*}
%%%END
%%%BLOCK id=262-03 ch=04 sec=竖直面圆周·绳模型 kind=derivation alt=none cont=none y=0.46-0.67
不满足 $V\ge\sqrt{gL}$ 时，
%FIG p262_b 0.040 0.495 0.274 0.660
\notefig[0.3]{figures/p262_b.png}
图中标注：做平抛；$V$ 损失$=0$
\[ V_0=\sqrt{\frac{gL}{2}}<\sqrt{gL} \]
\begin{align*}
V_0^2t^2&=L^2\sin^2\alpha\\
\frac{1}{2}gt^2&=L(1+\cos\alpha)
\end{align*}
$\Rightarrow\ \alpha=90^\circ$\ 圆心等高处\ 绷直
\[ mgl=\frac{1}{2}(V^2-V_y^2)\qquad V^2=V_y^2+2gl=4gl \] % CHECK: 「½」与括号之间似缺 m（原稿有一小撇）
\[ F-mg=\frac{mV^2}{L}\ \Rightarrow\ F=5mg \]
%%%END
%%%BLOCK id=262-04 ch=04 sec=竖直面圆周·等效重力加速度 kind=derivation alt=none cont=none y=0.69-0.92
\textbf{等效重力加速度}
%FIG p262_c 0.010 0.705 0.255 0.815
\notefig[0.3]{figures/p262_c.png}
\[ \left\{\begin{aligned}F+mg\sin\alpha&=\frac{mV_{\text{上}}^2}{L}\\ F&\ge0\end{aligned}\right. \qquad \begin{aligned}V_{\text{上}}&\ge\sqrt{gL\sin\alpha}\\ V_{\text{下}}&\ge\sqrt{5gL\sin\alpha}\end{aligned} \]
\[ mg\cdot2L\sin\alpha=\frac{1}{2}m(V_{\text{下}}^2-V_{\text{上}}^2)\qquad V_{\text{上}}^2=V_{\text{下}}^2-4gL\sin\alpha>gL\sin\alpha \] % CHECK: 末尾符号写作「>」，其下另起一行写 gLsinα；按前文类比应为「≥」
\[ \Delta F_{\text{上下}}=6mg\sin\alpha \]
%%%END
%%%BLOCK id=263-01 ch=04 sec=竖直面圆周运动·等效重力场 kind=method alt=09 cont=prev y=0.00-0.31
%FIG p263_a 0.06 0.005 0.18 0.115
\notefig[0.25]{figures/p263_a.png}
$\downarrow$ $E=\dfrac{2mg}{q}$
\[
\left\{\begin{aligned}
&mg+qE+F=\dfrac{mv^2}{L}\\
&T\ge 0\qquad V_{\text{上}}\ge\sqrt{3gL}
\end{aligned}\right.
\]
% CHECK: 第一式写 F、条件写 T≥0，应为同一个量（绳的拉力）
\[ G'=3mg,\qquad g'=3g,\qquad V_{\text{下}}\ge\sqrt{15gL} \]
% CHECK: g'=3g 中“3”前有涂改痕迹

（图注：等效最高点）
%FIG p263_b 0.03 0.108 0.205 0.215
\notefig[0.3]{figures/p263_b.png}
\[ \sqrt2\,mg+F=\dfrac{mv^2}{L}\qquad F\ge 0 \]
\[ V\ge\sqrt{\sqrt2\,gL} \]
%FIG p263_c 0.09 0.22 0.22 0.315
\notefig[0.25]{figures/p263_c.png}
\[ mg\left(L+\dfrac{\sqrt2}{2}L\right)+qE\,\dfrac{\sqrt2}{2}L=\dfrac12m\left(V_0^2-V^2\right) \]
%%%END
%%%BLOCK id=263-02 ch=04 sec=水平面圆周运动·转弯 kind=method alt=02 cont=none y=0.32-0.59
\textbf{水平面圆周运动}\quad 必有侧\textsuperscript{向外}向滑动趋势
% CHECK: “向外”为行上方的插入小字，插入位置约在“侧”“向”之间，语义待核（可能想表达“必有向外侧滑的趋势”）

\textbf{汽车转弯}\quad 水平路面上：
%FIG p263_d 0.17 0.388 0.285 0.466
\notefig[0.25]{figures/p263_d.png}
\[ f=\dfrac{mv^2}{r}\le\mu mg\qquad v\le\sqrt{\mu gR} \]

\textbf{倾斜路面上}\quad 无侧滑趋势\quad $V_0$.
%FIG p263_e 0.335 0.485 0.495 0.59
\notefig[0.3]{figures/p263_e.png}
\begin{align*}
N\sin\alpha&=\dfrac{mv_0^2}{R}\\
N\cos\alpha&=mg
\end{align*}
\[ V_0=\sqrt{gR\tan\alpha}\qquad\text{与 }\mu\text{ 无关} \]
%%%END
%%%BLOCK id=263-03 ch=04 sec=竖直面圆周运动·杆模型 kind=method alt=none cont=none y=0.60-0.78
\textbf{杆模型}
%FIG p263_f 0.06 0.635 0.285 0.72
\notefig[0.3]{figures/p263_f.png}
以向下为正
\[ mg+F=\dfrac{mv^2}{R}\qquad F\text{向下} \]
\[ mg-(-F)=\dfrac{mv^2}{R}\qquad F\text{向上} \]
\[
\left\{\begin{aligned}
&0\le v<\sqrt{gR}\quad\text{推力}\quad v\uparrow\ \text{推力}\downarrow\\
&v>\sqrt{gR}\quad\ \ v\uparrow\ F_{\text{拉}}\downarrow
\end{aligned}\right.
\]
% CHECK: v>√(gR) 时 F拉 随 v 增大应增大，原稿箭头画的是 ↓
%FIG p263_g 0.70 0.615 0.885 0.72
\notefig[0.3]{figures/p263_g.png}
\[ F=\dfrac{mv^2}{R}-mg \]
%%%END
%%%BLOCK id=263-04 ch=04 sec=竖直面圆周运动·绳模型临界 kind=method alt=none cont=none y=0.80-0.92
\textbf{绳不松弛}\quad $N=0$ 临界.
%FIG p263_h 0.07 0.86 0.155 0.92
\notefig[0.25]{figures/p263_h.png}
\[ V_0\le\sqrt{2gL} \]
\[ V_0\ge\sqrt{5gL} \]
%%%END
%%%BLOCK id=263-05 ch=04 sec=竖直面圆周运动·脱离临界 kind=method alt=06 cont=none y=0.92-1.00
%FIG p263_i 0.09 0.93 0.285 1.0
\notefig[0.3]{figures/p263_i.png}
\[ V_0=\sqrt{\dfrac{7}{2}gR} \]
\begin{align*}
mg\sin\alpha&=\dfrac{mv^2}{R}\\
-mgR(1+\sin\alpha)&=\dfrac12m\left(gR\sin\alpha-\dfrac{7}{2}gR\right)
\end{align*}
% CHECK: 第二式 mg 后的 R 与左括号重叠，按能量关系取 R
\[ \Rightarrow\ \alpha=30^\circ \]
%%%END
%%%BLOCK id=264-01 ch=04 sec=竖直面圆周运动·综合例题 kind=problem alt=none cont=none y=0.02-0.34
\begin{problem}[例]
%FIG p264_a 0.15 0.02 0.28 0.115
\notefig[0.25]{figures/p264_a.png}
\begin{solution}
%FIG p264_b 0.125 0.14 0.25 0.215
\notefig[0.25]{figures/p264_b.png}
\begin{align*}
&T+mg\cos\theta=\dfrac{mv^2}{R}\\
&mg(1-\cos\theta)=\dfrac12mv^2-\dfrac12m(gR)
\end{align*}
\[ T=3mg(1-\cos\theta) \]
对人\ $N=Mg-T\cos\theta=Mg-3mg(1-\cos\theta)\cos\theta$

$\ge Mg-3mg\cdot\dfrac14=Mg-\dfrac34mg$

$T_{\max}-mg=\dfrac{mV_{\text{下}}^2}{R}$\qquad $V_{\text{下}}=\sqrt{5gR}$

对人、$N=Mg+\dfrac{mv^2}{R}=Mg+5mg$

故人对地压力 $N\in\left[Mg-\dfrac34mg,\ Mg+5mg\right]$
% CHECK: 两处“5mg”的“5”字形有涂改/重叠（第二处尤其明显，可能是 6）；按 T_max=mg+5mg=6mg 推算，N 或应为 Mg+6mg
\end{solution}
\end{problem}
%%%END
%%%BLOCK id=264-02 ch=04 sec=竖直面圆周运动·脱离临界 kind=method alt=none cont=none y=0.34-0.47
是否可以沿圆弧到 $N$ 点？
%FIG p264_c 0.10 0.358 0.34 0.412
\notefig[0.4]{figures/p264_c.png}
\begin{itemize}
\item $V_0=\sqrt{gR}$ 时平抛
\item $V_0=0$ 时
\end{itemize}
\[
\left\{\begin{aligned}
&mg\sin\alpha=\dfrac{mv^2}{R}\\
&mg(1-\sin\alpha)\,\unclear{R}=\dfrac12mv^2
\end{aligned}\right.
\qquad \sin\alpha=\dfrac23
\]
% CHECK: 第二式 ")" 后的符号像 F，应为 R（mgR(1-sinα)=½mv²）
\unclear{处}脱离.
%%%END
%%%BLOCK id=264-03 ch=04 sec=水平面圆周运动·双绳 kind=method alt=none cont=none y=0.49-1.00
\textbf{双绳竖直平面内做水平圆周运动}
%FIG p264_d 0.07 0.526 0.22 0.64
\notefig[0.25]{figures/p264_d.png}
① OA绷紧 OB恰好松弛无张力
\[ mg\tan37^\circ=m\omega_1^2d\qquad d=L\sin37^\circ\qquad \omega_1=\sqrt{\dfrac{5g}{4L}} \]
② OB绷紧 OA恰好伸直无张力
\[ mg\tan\unclear{53}^\circ=m\omega_2^2d\qquad \omega_2=\sqrt{\dfrac{20g}{9L}} \]
% CHECK: 手写像 tan57°，按图中 53° 及 ω_2=√(20g/9L) 取 53
OA紧 OB松 $\xrightarrow{\ \omega_1\ }$ OA OB都紧 $\xrightarrow{\ \omega_2\ }$ OB紧 OA松.（图中以弧线分段，分界点标 $\omega_1$、$\omega_2$）

%FIG p264_e 0.04 0.745 0.175 0.86
\notefig[0.25]{figures/p264_e.png}
\[ mg\tan\alpha=m\omega_1^2d\qquad \omega_1=\sqrt{\dfrac{3g}{4d}} \]
OA紧 OB松 $\xrightarrow{\ \omega_1\ }$ $F_A$ 不变，$F_B\uparrow$

%FIG p264_f 0.02 0.865 0.19 0.995
\notefig[0.25]{figures/p264_f.png}
\[
\left\{\begin{aligned}
&T_1\sin\alpha+T_2\sin\beta=m\omega^2d\\
&T_1\cos\alpha=T_2\cos\beta+mg
\end{aligned}\right.
\qquad \omega\uparrow\quad T_1\uparrow\quad T_2\uparrow
\]
%%%END
%%%BLOCK id=265-01 ch=04 sec=平抛运动·斜面问题 kind=method alt=none cont=none y=0.00-0.09
\textbf{垂直落至斜面}
%FIG p265_a 0.00 0.00 0.21 0.085
\notefig[0.25]{figures/p265_a.png}
\begin{align*}
&V_x=V_0\\
&V_y=gt\\
&V=\sqrt{V_x^2+V_y^2}
\end{align*}
\[ \tan\theta=\dfrac{V_0}{V_y}=\dfrac{V_0}{gt}\qquad t=\dfrac{V_0}{g\tan\theta} \]
%%%END
%%%BLOCK id=265-02 ch=04 sec=平抛运动·斜面问题 kind=method alt=none cont=none y=0.09-0.21
\textbf{切入斜面}
%FIG p265_b 0.00 0.085 0.19 0.145
\notefig[0.25]{figures/p265_b.png}
%FIG p265_c 0.085 0.12 0.27 0.215
\notefig[0.3]{figures/p265_c.png}
\begin{align*}
&t=\dfrac{V_1}{a_1}\qquad V_1^2=2a_1d\qquad t=\dfrac{V_0\tan\theta}{g}\qquad d=\dfrac{V_0^2\sin\theta\tan\theta}{2g}\\
&V_y=gt=\dfrac{V_0\tan\theta}{g}
\end{align*}
% CHECK: V_y=gt=V_0 tanθ/g 末尾分母多一个 g（应为 V_y=gt=V_0 tanθ）；t=V_1/a_1 中的 V_1、a_1 字迹偏小
%%%END
%%%BLOCK id=265-03 ch=04 sec=平抛运动·斜面问题 kind=method alt=none cont=none y=0.21-0.33
\textbf{抛上斜面}
%FIG p265_d 0.00 0.21 0.26 0.285
\notefig[0.3]{figures/p265_d.png}
\[
\left\{\begin{aligned}
x&=V_0t\\
y&=\dfrac12gt^2\\
x_{\text{合}}&=\sqrt{x^2+y^2}
\end{aligned}\right.
\qquad
\begin{aligned}
\tan\theta&=\dfrac{y}{x}=\dfrac{gt}{2V_0}\\
t&=\dfrac{2V_0\tan\theta}{g}
\end{aligned}
\]
%%%END
%%%BLOCK id=265-04 ch=04 sec=平抛运动·抛入圆弧面 kind=method alt=none cont=none y=0.33-0.47
\textbf{抛入圆弧面}
%FIG p265_e 0.03 0.358 0.238 0.465
\notefig[0.3]{figures/p265_e.png}
\[
\left\{\begin{aligned}
h&=\dfrac12gt^2\\
R\pm\sqrt{R^2-h^2}&=V_0t
\end{aligned}\right.
\]
%%%END
%%%BLOCK id=265-05 ch=04 sec=平抛运动·圆心角与速偏角 kind=method alt=none cont=none y=0.49-0.74
%FIG p265_f 0.05 0.49 0.33 0.60
\notefig[0.4]{figures/p265_f.png}
\textbf{圆心角 $=$ 速偏角}
%FIG p265_g 0.09 0.595 0.29 0.745 soft
\notefig[0.3]{figures/p265_g.png}
%%%END
%%%BLOCK id=266-04 ch=04 sec=竖直面圆周运动·脱离临界 kind=derivation alt=06 cont=none y=0.82-1.00
%FIG p266_d 0.04 0.81 0.215 0.905
\notefig[0.3]{figures/p266_d.png}
\[
\left\{\begin{aligned}
&-mgR\sin\alpha=\dfrac12m\left(v^2-v_0^2\right)\\
&mg\sin\alpha=\dfrac{mv^2}{\unclear{R}}\\
&v\sin\alpha\,t=R\cos\alpha\\
&R\sin\alpha=v\cos\alpha\,t-\dfrac12gt^2
\end{aligned}\right.
\]
% CHECK: 第二式分母手写像 v，应为 R；最后一行紧贴页面底边，可能被裁切，其后或还有内容
%%%END
%%%BLOCK id=267-01 ch=04 sec=竖直面圆周运动·等效重力场 kind=problem alt=09 cont=none y=0.00-0.25
\begin{problem}
%FIG p267_a 0.07 0.02 0.36 0.18
\notefig[0.4]{figures/p267_a.png}
（图注：等效最低点）\quad $E=\dfrac{mg}{q}$\qquad 求 $F_{\max}$
\begin{solution}
\[
\left\{\begin{aligned}
&\sqrt2\,mg\sqrt2\,R=\dfrac12mv^2\\
&F-\sqrt2\,mg=\dfrac{mv^2}{R}\\
&\sqrt2\,mgR\left(1-\dfrac{\sqrt2}{2}\right)=\dfrac12m\left(v'^2-v_x^2\right)
\end{aligned}\right.
\]
% CHECK: 第三式括号内 v'^2 的撇号/指数手写不清
\end{solution}
\end{problem}
%%%END
%%%BLOCK id=269-04 ch=04 sec=圆周运动·圆锥摆 kind=method alt=none cont=none y=0.73-0.87
%FIG p269_c 0.64 0.725 0.86 0.87
\notefig[0.25]{figures/p269_c.png}
等效绳长$L$：
\[
\cos\alpha=\frac{g}{L\omega^2}
\]
其下方画有箭头$\downarrow$，再下方写$\alpha\uparrow$。图中转轴处标$\omega'$（带转向弧线），竖直虚线为转轴，斜绳与竖直方向夹角处标$\alpha$。
%%%END
%%%BLOCK id=269-06 ch=04 sec=圆周运动·离心现象 kind=note alt=05 cont=none y=0.93-0.96
太空中圆周问\unclear{题}\quad $V$越大\ 油水越易分离
%%%END
%%%BLOCK id=270-02 ch=04 sec=平抛运动·空间极值 kind=note alt=none cont=none y=0.09-0.12
平抛在空间中对角线处最远
%%%END
%%%BLOCK id=271-01 ch=04 sec=水平面圆周·摩擦力提供向心力 kind=problem alt=02 cont=none y=0.00-0.19
% 贴纸印刷题（上缘被截断）；原稿上有铅笔划线：“力支撑冰面才能保证”处有划线，“过弯道滑行时的运动员”处有划线
\begin{problem}[贴纸印刷题]
\unclear{…}2022年\unclear{…}威、曲春雨、范可新、武大靖和张雨婷组成的中国队夺得北京冬奥会短道速滑男女2000米混合接力冠军，为中国体育代表团收获了北京冬奥会的首枚金牌。短道速滑运动员过水平弯道时常用手支撑冰面以防侧滑，某运动员质量为$75\unit{kg}$，某次过弯道时的半径为$25\unit{m}$，速率为$36\unit{km/h}$，冰刀与冰面间的动摩擦因数为$0.2$，手套与冰面间的动摩擦因数为$0.8$，重力加速度$g=10\unit{m/s^2}$。过弯道滑行时的运动员手脚距离相对半径可忽略，弯道滑行的过程视为一段圆周运动，则该运动员至少用多大的力支撑冰面才能保证不发生侧滑（\struck{B}C） % 括号内手写答案：“B”被斜线划去，旁写C

A.\ $150\unit{N}$\qquad B.\ $200\unit{N}$\qquad C.\ $250\unit{N}$\qquad D.\ $300\unit{N}$

手写铅笔批注（部分被印刷图片遮挡，数字难辨）：
\begin{itemize}
\item 第2行右端圈出：$10\unit{m/s}$
\item $0.2+0.8(\unclear{780}+N)=75\cdot\dfrac{\unclear{100}}{\unclear{30}}$，\ 其下一行写$0.8(\unclear{780}+N)$（括号部分划有下划线），再下一行$=300$ % CHECK: “780”疑为750；式子本身（0.2、0.8的组合）按原稿照录
\item 印刷图片上写有：$0.6$
\item 右缘：$7.5$（下有横线）；圈出的$300$；其下$\unclear{1000}$（被页边截断）
\end{itemize}
\end{problem}
%%%END
%%%BLOCK id=272-01 ch=04 sec=竖直面圆周·斜抛综合 kind=problem alt=06 cont=none y=0.03-0.30
%FIG p272_a 0.04 0.03 0.42 0.13 soft
\notefig[0.5]{figures/p272_a.png}
图中标注：$m=0.02\unit{kg}$，$V_0$（箭头），$R=0.4\unit{m}$，粗糙段$\mu=0.25$，斜坡端有$\alpha$。
\begin{problem}
(1) 滑块始终在轨道上运动，求$V_0$范围

(2) 滑块恰过D，求$E_{k\min}$时$X$与$\alpha$的关系，并求$X$最大值。
\begin{solution}
从\unclear{…}到A时速度为0
\[
mg(h-R)-\mu mgR=0-0\qquad h=0.5\unit{m}
\]
\[
mg(R-h)-\mu mgR=0-\frac{1}{2}mV_{01}^2 % 原稿“μ”写得像“4”
\]
\[
\therefore V_{01}=2\unit{m/s}\qquad 0\le V_{01}\le2\unit{m/s}
\]
\[
mg\,2R\cos\alpha=\frac{1}{2}mV_F^2
\]
\[
X=V_F\cos\alpha\cdot\frac{V_F\sin\alpha}{g}=\frac{V_F^2}{g}\sin\alpha\cos\alpha=1.6\cos^2\alpha\sin\alpha\le\frac{16\sqrt3}{45}\ \Rightarrow\ \tan\alpha=\frac{\sqrt2}{2}\ \text{时取等}
\]
\end{solution}
\end{problem}
%%%END
%%%BLOCK id=277-02 ch=04 sec=类斜抛·恒力 kind=problem alt=03 cont=none y=0.42-0.79
% 图：直线ab，a处速度 V_a=3V(水平向右)，b处速度 V_b=4V；图中另标 \alpha、53^\circ、\theta、\beta、37^\circ，以及竖直向下的“设F风”
%FIG p277_c 0.055 0.425 0.28 0.61
\notefig[0.3]{figures/p277_c.png}

\[
V_a\sin(\alpha+\theta)=V_b\sin(\theta-\beta)
\]
$\therefore\theta=74^\circ$

从$a\to b$ 平均 $\bar V=\dfrac{3V\cos53^\circ+4V\cos37^\circ}{2}=2.5V$

$t=\dfrac{d}{V}=\dfrac{2d}{5V}$ % CHECK: 手写第一个分母无上横线，按上下文应为平均速度 \bar V

$a\to b$\quad $a=\dfrac{V_b\cos37^\circ-(-V_a\cos53^\circ)}{t}=\dfrac{25V^2}{2d}$ % CHECK: 分子“25”手写偏草，按推算为 25

$a\to b$ 风力 $F=ma=\dfrac{25mV^2}{2d}$

小球沿风力方向速度最小时 $V$最小 % 原稿“球”字写成“王求”

$V_{\min}=V_a\sin53^\circ=2.4V$
%%%END
%%%BLOCK id=285-11 ch=04 sec=平抛运动·竖直方向匀加速 kind=note exp=1 alt=01 cont=none y=0.62-0.79
平抛的竖直方向为匀加\unclear{速}

%FIG p285_d 0.10 0.665 0.27 0.785
\notefig[0.25]{figures/p285_d.png}

\begin{gather*}
h_1-h_2=h_2-h_3\\
h_2-h_3=h_3-h_4
\end{gather*}
\[ h_1+3H_3=3H_2+h_4 \]% CHECK: 手写 H_3、H_2 为大写，应为 h_3、h_2
%%%END
%%%BLOCK id=285-12 ch=04 sec=平抛运动·轨迹与曲杆 kind=note alt=none cont=none y=0.82-0.88
% 左下角被折页遮挡，“后 E⃗”之后到“y=”之前的内容被遮住
后 $\vec{E}$\unclear{…}\quad $y=\dfrac{gx^2}{2v_0^2}$ 当符合曲杆轨迹时，$N=0$，是平抛。$mg=\dfrac{mv_0^2}{\rho}$ 曲率半径。
%%%END
%%%BLOCK id=286-05 ch=04 sec=竖直面圆周·临界 kind=note alt=06 cont=none y=0.85-0.96
% 图中标注：顶部小球 V1=√(gr)，底部小球 V2=√(5gr)
%FIG p286_d 0.04 0.855 0.215 0.96
\notefig[0.25]{figures/p286_d.png}

不可能同时位于水平直径两端

\unclear{系统 $E_{p\text{低}}>2mgr$}\quad 不会碰撞
%%%END
%%%BLOCK id=289-01 ch=04 sec=斜面上的圆周运动·重力功率 kind=note alt=06 cont=none y=0.00-0.13
%FIG p289_a 0.095 0.0 0.43 0.125
\notefig[0.3]{figures/p289_a.png}

$P_G$ 先↑后↓

$P_{G\max}$ 时 $F_{\text{合}}$ 与 $OA$ 平行
%%%END
%%%BLOCK id=289-02 ch=04 sec=运动的合成与分解·关联速度 kind=problem alt=none cont=none y=0.13-0.27
%FIG p289_b 0.105 0.135 0.29 0.26
\notefig[0.3]{figures/p289_b.png}

%FIG p289_c 0.325 0.115 0.65 0.258
\notefig[0.4]{figures/p289_c.png}

$V_A\sin\theta=V_B\cos\dfrac{\theta}{2}$

两球速度大小相等时 $B$ 在 $O$ 点正上方% CHECK: “O点”手写像 O（图中对应点标注像 D/O）
%%%END
%%%BLOCK id=294-01 ch=04 sec=平抛·落在曲面上 kind=problem alt=06 cont=none y=0.03-0.40
\begin{problem}[1]
%FIG p294_a 0.07 0.06 0.27 0.15
\notefig[0.3]{figures/p294_a.png}

\begin{solution}
$v_0$ 越小 $h$ 越高 飞行时间越长
\begin{align*}
R\cos\theta&=v_0t\\
R\sin\theta&=\frac{1}{2}gt^2\\
mgR\sin\theta&=E_K-\frac{1}{2}mv_0^2\\
E_K&=mgR\left(\frac{3}{4}\sin\theta+\frac{1}{4\sin\theta}\right)\\
&>\frac{\sqrt{3}}{2}mgR % CHECK: 原稿为“>”，应为 $\ge$
\end{align*}
$P_{\min}=\sqrt{2mE_{k\min}}$\quad 此时 $x=\dfrac{\sqrt{6}}{3}R$\quad $v_0=\sqrt{\dfrac{\sqrt{3}}{3}gR}$
\end{solution}
\end{problem}
%%%END
